%% Manuscript: Site Effects in Multi-Centre EEG Data for Neuroscreening
%% Target journal: Biomedical Signal Processing and Control (Elsevier)
%% LaTeX class: elsarticle, numerical citations
%%
%% Compile with:
%%   pdflatex manuscript_eeg_harmonization.tex
%%   bibtex   manuscript_eeg_harmonization
%%   pdflatex manuscript_eeg_harmonization.tex
%%   pdflatex manuscript_eeg_harmonization.tex

\documentclass[preprint,12pt]{elsarticle}

%\documentclass[final,3p,times]{elsarticle}

\usepackage{amsmath}
\usepackage{amssymb}
\usepackage{graphicx}
\usepackage{booktabs}
\usepackage{multirow}
\usepackage{hyperref}
\usepackage{lineno}
\usepackage{microtype}
\usepackage{xcolor}
\usepackage{subcaption}
\usepackage{setspace}

\newcommand{\td}[1]{\textcolor{orange}{\textbf{[TODO: #1]}}}

%\linenumbers
%\doublespacing

\journal{Biomedical Signal Processing and Control}

% -----------------------------------------------------------------------
\begin{document}

\begin{frontmatter}

\title{Harmonisation of Multi-Site EEG Data for Clinical Neuroscreening:
       Quantifying the Site Effect and Comparing Correction Strategies}

\author[aff1]{Kacper Chorzela}
\author[aff1]{Julia Pelczar}
\author[aff1]{Jaros\l{}aw {\.Z}ygierewicz\fnref{fn1}\corref{cor1}}
\ead{jarekz@fuw.edu.pl}

\cortext[cor1]{Corresponding author.}
\fntext[fn1]{ORCID: 0000-0002-7536-0735}

\affiliation[aff1]{organization={Faculty of Physics, University of Warsaw},
             addressline={Pasteura 5},
             city={Warsaw},
             postcode={02-093},
             country={Poland}}

% -----------------------------------------------------------------------
\begin{abstract}
Pooling electroencephalographic (EEG) recordings from multiple clinical
centres introduces systematic, non-biological variance---the
\emph{site effect}---that can severely degrade the generalisation of
machine-learning diagnostic models.
We characterise the site effect in a large-scale corpus of 54\,779
clinical EEG recordings from 30 Polish hospitals
($\mathrm{ELM}_s$) and evaluate a comprehensive range of
harmonisation strategies under a unified experimental design.

Feature-level harmonisation methods---site-wise standardisation,
neuroCombat, CovBat, and PCA-based feature-space
pre-orthogonalisation---are assessed on the normal-recording subset
($\mathrm{ELM}_n$).
Signal-level harmonisation via a Domain Adversarial Neural Network
(DANN) integrated with the multiple-instance learning architecture
of~\cite{poziomska2026quantity} is evaluated as a complementary
end-to-end approach.
Throughout, a \emph{two-classifier evaluation framework}---a logistic
regression linear probe and a CatBoost nonlinear classifier applied
symmetrically to every condition---separates the linearly accessible
component of the site effect from nonlinear inter-feature covariance
structure.
Clinical performance is assessed exclusively under a rigorous
Leave-One-Site-Out (LOSO) protocol, directly simulating zero-shot
deployment to a previously unseen recording environment.

A CatBoost classifier trained on the unharmonised 2\,850-dimensional
feature set (Riemannian covariance projections, multitaper power
spectra, and inter-electrode coherence) achieved a Matthews
Correlation Coefficient (MCC) of $0.843 \pm 0.007$, confirming that
site identity is strongly encoded in the feature space.
SHAP analysis identified inter-electrode coherence---particularly
O1--O2 and beta-band temporal coherence---as the primary carriers
of site-specific information.

A central finding is the \emph{harmonisation paradox}: all marginal
correction methods reduced inter-site Cohen's~$d$ and drove the
linear probe to near-chance MCC, yet simultaneously increased the
nonlinear CatBoost MCC towards a near-ceiling value, demonstrating
that feature-by-feature correction reorganises rather than eliminates
multivariate covariance-based site structure.
PCA pre-orthogonalisation substantially attenuates this paradox and
yields more uniform cross-site neuroscreening performance.

\textbf{[PLACEHOLDER: Complete aggregate results for all conditions
once all DANN experiments are finalised.]}
\end{abstract}

\begin{keyword}
EEG \sep site effect \sep batch effect \sep data harmonisation \sep
ComBat \sep CovBat \sep PCA pre-orthogonalisation \sep neuroscreening
\sep multi-centre study \sep machine learning \sep SHAP \sep DANN \sep
domain adaptation \sep logistic regression \sep two-classifier framework
\end{keyword}

\end{frontmatter}

\td{
\texttt{
TO GENERATE --- Priority 1 (required for submission):\\
\\
\quad [Table 3] Two-classifier site-classification MCC for all conditions\\
\quad\quad (Tab.~tab:site\_class\_harmonisation):\\
\quad\quad -- LR-MCC and CatBoost-MCC for: CovBat, all PCA conditions,\\
\quad\quad\quad MINet, MINet-MTL, MINet-DANN\\
\quad\quad -- LR-MCC for: Unharmonised, Site-wise, neuroCombat\\
\\
\quad [Table 4] Aggregate LOSO neuroscreening MCC and AUC\\
\quad\quad (Tab.~tab:clinical\_loso):\\
\quad\quad -- CovBat, all PCA conditions, MINet-DANN-CF\\
\\
\quad [Figure A] LR vs CatBoost MCC per harmonisation condition\\
\quad\quad (fig:lr\_cb\_gap) -- grouped paired bar chart, three blocks:\\
\quad\quad -- Block 1: feature-level without PCA\\
\quad\quad -- Block 2: feature-level with PCA\\
\quad\quad -- Block 3: signal-level (on h)\\
\quad\quad -- whiskers = SD; annotation of Delta MCC per pair\\
\quad\quad -- replaces Figs site\_mcc\_pca\_00/10 in main body\\
\quad\quad\quad (those move to Appendix)\\
\\
\quad [Figure B] Domain MCC per signal-level variant\\
\quad\quad (fig:dann\_domain\_mcc) -- currently commented out;\\
\quad\quad -- bar chart: MINet / MINet-MTL / MINet-DANN (/ MINet-DANN-CF)\\
\quad\quad -- to be placed in subsec:harmonisation\_results\\
\quad\quad\quad alongside fig:lr\_cb\_gap\\
\\
\quad [Figure C+D] Per-institution LOSO AUC and MCC\\
\quad\quad for all signal-level conditions including MINet-MTL\\
\quad\quad (extend figs all\_strategies\_auc/mcc or add separate panels)\\
\quad\quad -- MINet-MTL per-institution profile currently missing;\\
\quad\quad\quad high SD=0.120 warrants visualisation\\
\quad\quad -- MINet-DANN-CF per-institution profile (TBF)\\
\\
TO GENERATE --- Priority 2:\\
\\
\quad [Appendix] Per-institution site-classification MCC\\
\quad\quad without PCA (fig:site\_mcc\_no\_pca) -- move from main body\\
\quad\quad with PCA (fig:site\_mcc\_pca) -- move from main body\\
\\
\quad [Appendix] Per-institution LOSO AUC\\
\quad\quad without PCA (fig:patho\_auc\_no\_pca) -- currently in main body,\\
\quad\quad\quad consider moving if length is an issue\\
\quad\quad with PCA (fig:patho\_auc\_pca) -- idem\\
\\
PENDING AUTHOR DECISIONS:\\
\quad -- full-data checkpoint identity in MINet LOSO paragraph\\
\quad -- epoch selection criterion for MINet-MTL in LOSO\\
\quad -- hyperparameter source for CatBoost neuroscreening classifier\\
\quad\quad (subsubsec:classifier, current \td{sprawdzic})\\
}}





% -----------------------------------------------------------------------
\section{Introduction}
\label{sec:introduction}

Electroencephalography (EEG) is a non-invasive, cost-effective
neurophysiological technique that holds substantial promise for
scalable, data-driven clinical diagnostics~\cite{poziomska2026quantity}.
The increasing availability of large, routinely collected clinical EEG
archives has opened new avenues for training machine-learning (ML)
classifiers capable of detecting neurological abnormalities at a
population scale~\cite{mushtaq2025eeg}.
A prerequisite for such systems to be deployable in practice is the
ability to generalise across acquisition environments: a model trained
on recordings from one set of hospitals must perform reliably on data
from institutions not represented during training.

This requirement is fundamentally challenged by the \emph{site
effect}---systematic, non-biological variance arising from
inter-centre differences in EEG hardware, electrode placement
conventions, amplifier settings, recording protocols, patient
preparation procedures, and patient
demography~\cite{jaramillo2024combat,li2022harmnqeeg}.
When ML models are trained on multi-centre data without addressing
the site effect, they can exploit institution-specific signatures in
the feature space rather than learning clinically discriminative
patterns, leading to marked degradation of performance when the model
encounters data from a previously unseen
hospital~\cite{kiessner2023extended,mellot2023meeg}.
As large, heterogeneous clinical EEG corpora become the natural
substrate for the next generation of diagnostic algorithms, the site
effect has emerged as one of the principal obstacles to building
generalisable neuroscreening models~\cite{dong2025normative,
mushtaq2025eeg}.

Statistical harmonisation methods, most prominently the ComBat
family~\cite{johnson2007adjusting,fortin2017harmonization,
fortin2018cortical}, were developed to address the analogous
batch-effect problem in genomics and neuroimaging, and have more
recently been applied to EEG
features~\cite{jaramillo2024combat,li2022harmnqeeg}.
CovBat~\cite{hu2023image} extends this framework to the multivariate
setting by additionally harmonising batch-specific structure in the
principal-component space of the residual covariance matrix.
Although computationally lightweight and widely available, all such
methods share a fundamental constraint: they operate feature by feature
and may therefore leave intact site-related structure encoded in
inter-feature covariances.

An alternative paradigm embeds harmonisation directly into the
representation-learning objective via Domain Adversarial Neural
Networks (DANNs)~\cite{ganin2016domain}, which have demonstrated
improved cross-subject and cross-site generalisation on EEG
classification tasks~\cite{bethge2022multidomain,liu2023multisource,
ozdenizci2020invariant}.
Critically, however, whether DANN-based harmonisation genuinely
suppresses the \emph{nonlinear} covariance component of the site
effect---or merely its linearly accessible component---has not been
systematically investigated.

The present study pursues two overarching goals.
First, we provide a rigorous characterisation of the site effect in a
corpus of 54\,779 EEG recordings from 30 Polish hospitals using a
\emph{two-classifier evaluation framework} that separates linearly
accessible site information from nonlinear inter-site covariance
structure.
Second, we evaluate a comprehensive range of harmonisation strategies
under a unified experimental design, spanning simple site-wise
standardisation, neuroCombat, CovBat, PCA pre-orthogonalisation, and
three signal-level adversarial variants, and critically assess the
conditions under which these methods succeed or fail under a rigorous
Leave-One-Site-Out (LOSO) protocol.

The main contributions of this paper are as follows.
\begin{enumerate}
  \item We demonstrate that the site effect in a large clinical EEG
        corpus is strong and broadly distributed, with a nonlinear
        classifier achieving $\mathrm{MCC} = 0.843 \pm 0.007$ in a
        site identification task.
  \item We identify, via SHAP analysis and ablation experiments, that
        inter-electrode coherence features are the primary carriers of
        site-specific information, with O1--O2 coherence and beta-band
        temporal coherence being particularly prominent.
  \item We introduce a \emph{two-classifier evaluation framework}
        that uses a logistic regression linear probe and a CatBoost
        nonlinear probe applied symmetrically to every harmonisation
        condition, separating linearly accessible site information
        from nonlinear covariance structure.
  \item We reveal and characterise a \emph{harmonisation paradox}:
        marginal correction methods reduce the linear site-effect
        component (Cohen's~$d$, linear probe MCC near chance) yet
        simultaneously increase the nonlinear site-classification
        MCC towards a near-ceiling value, because equalising marginal
        distributions amplifies the salience of the surviving
        covariance-based site signatures.
  \item We show that PCA pre-orthogonalisation substantially
        attenuates the harmonisation paradox and stabilises
        cross-site neuroscreening performance across institutions.
  \item We evaluate three signal-level training conditions—a non-adversarial baseline (trained from scratch), a multi-task network (optimizing both diagnosis and site prediction), and a full adversarial framework (DANN)—under the same Leave-One-Site-Out (LOSO) protocol. This allows us to compare end-to-end representation learning in terms of both site discriminability and downstream neuroscreening performance. Our results demonstrate that while the naive multi-task approach heavily encodes site-specific features (Domain MCC = 0.568±0.077), the adversarial DANN successfully creates a site-invariant representation (Domain MCC near zero: 0.009±0.014) without compromising high diagnostic accuracy (Target AUC = 0.810±0.073), effectively matching the non-adversarial baseline (Target AUC = 0.811±0.070). We additionally evaluate whether classifier recalibration on the frozen adversarial backbone (MINet-DANN-CF) further stabilizes the results. As hypothesized, this recalibration recovers and slightly improves the optimal diagnostic performance (Target AUC = 0.812±0.071, Target MCC = 0.475±0.130) without altering the underlying site-invariant features.
\end{enumerate}

% -----------------------------------------------------------------------
\section{Related Work}
\label{sec:related}

\subsection{The Site Effect in Multi-Centre EEG Studies}
\label{subsec:site_effect}

Systematic, non-biological variability arising from differences in
data acquisition conditions is a well-established problem across
quantitative biomedical sciences, where it is collectively termed the
\emph{batch effect}~\cite{johnson2007adjusting}.
In EEG research, the equivalent phenomenon is termed the \emph{site
effect}~\cite{li2022harmnqeeg}.
As large multi-centre EEG corpora become increasingly central to
data-driven clinical neuroscience~\cite{dong2025normative,
mushtaq2025eeg}, the site effect has emerged as a primary obstacle to
building classifiers that generalise beyond the hospitals represented
in training data.

Site effects arise from a heterogeneous combination of technical and
procedural factors including amplifier design, electrode materials,
impedance characteristics, recording-pipeline transfer functions,
activation procedures, resting-state recording duration, conventions
regarding eye-open versus eye-closed epochs, scalp preparation, and
patient alertness
management~\cite{jaramillo2024combat,mushtaq2025eeg}.
Demographic and clinical case-mix differences across institutions
further contribute to apparent site differences even when acquisition
protocols are nominally identical~\cite{dong2025normative}.
Because these sources act simultaneously and interact in complex ways,
the site effect is not a simple additive shift in individual features
but a high-dimensional, multivariate structure distributed across the
full feature space~\cite{mellot2023meeg}.

The degree to which site identity is encoded in EEG feature spaces
has been documented quantitatively.
Li~et~al.~\cite{li2022harmnqeeg} demonstrated systematic batch
effects across 12 devices and 9 countries in the HarMNqEEG corpus.
Jaramillo-Jimenez~et~al.~\cite{jaramillo2024combat} confirmed, using
$t$-SNE visualisation and Kruskal--Wallis tests, that spectral
features from five independent resting-state EEG datasets cluster by
site, with particularly prominent effects in the aperiodic offset and
oscillatory alpha power.
Kiessner~et~al.~\cite{kiessner2023extended} observed marked
performance degradation when a deep neural network trained on one
clinical corpus was evaluated on data from a different centre.
Mellot~et~al.~\cite{mellot2023meeg} showed that even
covariance-based EEG representations suffer substantial dataset shift
across acquisition devices and populations.

\subsection{Statistical Harmonisation Methods}
\label{subsec:statistical_harmonisation}

The canonical method is ComBat, originally proposed by
Johnson~et~al.~\cite{johnson2007adjusting} for batch-effect
correction in microarray gene-expression data.
Fortin~et~al.~\cite{fortin2017harmonization,fortin2018cortical}
established ComBat as the standard for retrospective neuroimaging
harmonisation.
CovBat~\cite{hu2023image} extends neuroCombat by additionally
harmonising the principal-component space of the residual covariance
matrix, thereby addressing batch effects residing in inter-feature
covariance structure.
CovBat has been shown to outperform standard ComBat in both
multivariate statistical evaluations and downstream machine-learning
metrics~\cite{hu2023image}.

In the EEG domain, Jaramillo-Jimenez~et~al.~\cite{jaramillo2024combat}
evaluated four ComBat variants on pooled resting-state EEG data from
five independent datasets and found that all variants reduced
batch-related dispersion in spectral features while preserving
age-related associations.
Li~et~al.~\cite{li2022harmnqeeg} demonstrated that ComBat-based
normalisation substantially reduced device-related clustering across
12 devices and 9 countries.

The most critical limitation of all ComBat-family methods is shared:
adjusting mean and variance independently for each feature leaves
intact any site-related structure residing in the \emph{covariance}
between features~\cite{hu2023image}.
Additionally, CovBat's principal-component decomposition step may be
unstable when individual sites have small sample
sizes~\cite{hu2023image}, and all ComBat-family methods require all
sites to be present at parameter-estimation time and cannot directly
accommodate recordings from a site not represented in the training
data~\cite{johnson2007adjusting}.
The present study provides the first systematic empirical evidence for
the covariance limitation in a large clinical EEG corpus through the
two-classifier evaluation framework.

\subsection{Deep Learning and Adversarial Harmonisation}
\label{subsec:dann_related}

An alternative paradigm embeds harmonisation directly into the
representation-learning objective, training a neural network to
produce features that are simultaneously predictive of the clinical
label and invariant to the recording site.
The canonical architecture is the Domain Adversarial Neural
Network~\cite{ganin2016domain}, which consists of a shared feature
extractor, a label predictor, and a domain discriminator connected via
a Gradient Reversal Layer (GRL).

Adversarial representation learning has been applied widely in the
EEG literature.
\"{O}zdenizci~et~al.~\cite{ozdenizci2020invariant} demonstrated
significant improvements in cross-subject generalisation on motor
imagery tasks.
Liu~et~al.~\cite{liu2023multisource} proposed MTDANN, which uses
EEGNet~\cite{lawhern2018eegnet} as the shared feature extractor and
outperformed single-source DANN and transfer learning baselines on
three public motor imagery datasets.
Bethge~et~al.~\cite{bethge2022multidomain} showed that a unified
adversarial invariance objective across multiple emotion recognition
datasets suppressed data-source-related information leakage by 35\%
while maintaining classification performance.
The adversarial approach developed in the present work
(Section~\ref{sec:approach2}) integrates the DANN objective
with the multiple-instance learning architecture of
Poziomska~et~al.~\cite{poziomska2026quantity}, which was designed
specifically for multi-recording EEG pathology decoding at scale.

An important open question, which the present study directly
addresses, is whether adversarial training suppresses the
\emph{nonlinear covariance} component of the site effect or only its
linearly accessible component.
The two-classifier framework---applying both a linear probe and a
CatBoost nonlinear probe to the DANN representations---is designed to
provide a direct empirical answer to this question
(Section~\ref{subsec:harmonisation_results}).

Despite its conceptual advantages, adversarial domain adaptation
introduces practical challenges that must be carefully managed.
The minimax training objective can be difficult to optimise: if the
domain classifier converges too quickly relative to the feature
extractor, the gradient signal it provides becomes uninformative,
while slow convergence leaves domain-specific structure in the
representation.
The adaptation weight $\lambda$ must therefore be scheduled carefully;
the progressive schedule of Ganin~et~al.~\cite{ganin2016domain} is
designed to avoid early instability by increasing $\lambda$ from 0 to
1 over the course of training.
In practice, unconstrained adversarial training from random
initialisation can lead to catastrophic forgetting: the model reaches
an optimal diagnostic checkpoint early in training but then degrades
as the adversarial objective distorts the decision boundaries.
Unlike ComBat, which produces an interpretable corrected feature set,
DANN-style approaches yield latent representations whose
site-invariance cannot be directly audited and must be assessed via
proxy metrics such as site-classification MCC.
These trade-offs motivate the comparative evaluation and the
classifier recalibration strategy described in
Section~\ref{sec:training_variants}.

% -----------------------------------------------------------------------
\section{Materials and Methods}
\label{sec:methods}

\subsection{Dataset}
\label{subsec:dataset}

\subsubsection{The ELM19 Corpus and the
  \texorpdfstring{$\mathrm{ELM}_s$}{ELMs} Subset}

We use the ELM19 database~\cite{poziomska2026quantity}, a large-scale
multi-centre EEG corpus provided by Elmiko Biosignals sp.~z~o.o.\
(Milanówek, Poland), comprising 55\,787 clinical EEG recordings from
39 medical institutions across Poland during routine clinical
diagnostics.
Each recording is annotated with a binary normal/pathological label
derived from the accompanying EEG report following the procedure
detailed in~\cite{poziomska2026quantity}.
Patient ages range from 16 to 70 years; 52.9\% of recordings are
from female patients; and 52.2\% of recordings are labelled as
pathological.
EEG signals were acquired using the international 10--20 electrode
placement system with 19 channels, with sampling rates varying
between 200 and 500\,Hz across institutions.

Sites contributing fewer than 50 normal recordings were excluded to
ensure sufficient within-site representation for harmonisation
parameter estimation.
The resulting dataset---hereafter $\mathrm{ELM}_s$---consists of
54\,779 EEG recordings from 30 medical institutions, comprising
26\,568 normal (48.5\%) and 28\,211 pathological (51.5\%) recordings
from 28\,960 female subjects (52.9\%).
The distributions of recordings across all 39 institutions before and
after this filtering step are illustrated in
Figure~\ref{fig:institution_histograms}.

\begin{figure}[htbp]
  \centering
  \begin{subfigure}[t]{0.48\textwidth}
    \centering
    \includegraphics[width=\textwidth]{%
      Figs/ELM19_excluded_institutions_normal_recordings.png}
    \caption{Number of normal recordings per institution.}
    \label{fig:normal_hist_excluded}
  \end{subfigure}
  \hfill
  \begin{subfigure}[t]{0.48\textwidth}
    \centering
    \includegraphics[width=\textwidth]{%
      Figs/ELM19_excluded_institutions_all_recordings.png}
    \caption{Total number of recordings (normal and pathological)
      per institution.}
    \label{fig:total_hist_excluded}
  \end{subfigure}
  \caption{Distribution of recordings across the 39 medical
    institutions in the ELM19 corpus.
    Institutions below the 50-normal-recording threshold (dashed line)
    were excluded from further analysis, yielding the
    $\mathrm{ELM}_s$ subset of 30 institutions.}
  \label{fig:institution_histograms}
\end{figure}

\subsubsection{Analysis Subset: Normal Recordings
  (\texorpdfstring{$\mathrm{ELM}_n$}{ELMn})}

To isolate the site effect from variance attributable to clinical
pathology, all primary analyses---exploratory data analysis, site
classification, harmonisation evaluation, and feature
importance---were conducted on the subset of recordings labelled as
normal, hereafter referred to as $\mathrm{ELM}_n$.
The full $\mathrm{ELM}_s$ dataset was used only in the downstream
neuroscreening experiment.

\subsubsection{Acquisition Metadata}
\label{subsubsec:acquisition_metadata}

All 30 institutions used EEG hardware and software manufactured by
Elmiko Biosignals sp.~z~o.o.
Of the 21 institutions for which hardware metadata were available,
17 used one amplifier model and 4 used a different model;
the remaining 9 institutions lacked hardware documentation
(Table~\ref{tab:acquisition_metadata}).
Software versions spanned a wide range, from early releases (v6.x)
through intermediate versions (v8.x, v10.x, v12.x) to recent builds
(v14.x, v15.x); several institutions changed software versions during
the data-collection period.
Sampling frequencies, as reported in EDF file headers, were
predominantly ${\sim}250$\,Hz, though minor variations existed
(e.g.,\ 256\,Hz at two sites and 300\,Hz at one site).
Acquisition-level filter settings recorded in the EDF headers also
differed across sites: 82.7\% of all recordings reported a 0\,Hz
high-pass setting (DC-coupled acquisition), while 16.5\% reported a
0.16\,Hz cutoff.
Low-pass acquisition filters similarly varied, with the majority of
institutions reporting 125\,Hz and a subset reporting lower values
(47--104\,Hz).
We note that EDF header metadata are written by the recording software
and may not always reflect the true analog hardware configuration;
the values reported here should therefore be interpreted as nominal
acquisition settings.
These metadata provide concrete evidence that the site effect has
identifiable technical component beyond purely demographic or clinical
differences.

\subsection{Inter-Institutional Heterogeneity}
\label{subsec:exploratory}

Exploratory analysis of the $\mathrm{ELM}_n$ subset revealed
substantial heterogeneity across the 30 institutions in patient
demographics and recording protocols.
Age distributions varied markedly between sites: institutions such as
H21, H29, and H30 served predominantly young populations (median
ages 19.5, 17.8, and 17.6 years, respectively), while institutions
such as H7 and H19 served predominantly older patients (median ages
50.8 and 52.1 years, respectively).
The maximum difference in median age between any two institutions
exceeded 30 years (full statistics are given in
Table~\ref{tab:age_stats} in the Appendix).
Sex distributions were broadly balanced across most institutions,
with minor imbalances at a small number of sites.
Recording duration showed further heterogeneity: while the majority
of centres produced sessions of approximately 1\,152 seconds
(median), institution H3 showed a bimodal distribution suggesting the
use of multiple recording protocols within a single institution, and
H16 recorded sessions lasting up to 8\,322 seconds.
Raw signal amplitudes (measured prior to preprocessing;
Section~\ref{subsec:preprocessing}), quantified as the median of per-channel median
absolute deviations (MAD) across the full recording, varied by more
than two orders of magnitude across institutions (median MAD range:
3.9--1\,611\,$\mu$V; Figure~\ref{fig:raw_amplitude}), likely
reflecting differences in amplifier gain and hardware high pass filters settings across hardware
configurations.
Frame rejection rates during preprocessing were generally low
(overall mean 1.9\%, median 0.3\%), with most institutions below
3\%; institution H22 was a notable outlier at 8.2\%, driven
primarily by flat-line electrode rejections.



Institution sample sizes in $\mathrm{ELM}_n$ ranged from 86 (H23)
to 6\,222 (H10), a 72-fold imbalance.
The proportion of normal recordings in $\mathrm{ELM}_s$ also varied
widely across institutions, from 7.6\% (H30) to 78.3\% (H7), further
motivating the restriction to normal recordings ($\mathrm{ELM}_n$)
for site-effect analyses.

%Together, these observations indicate that inter-site feature differences in $\mathrm{ELM}_n$ reflect a mixture of acquisition and population-level sources.


\subsection{EEG Preprocessing}
\label{subsec:preprocessing}

Raw EEG signals were preprocessed using
MNE-Python~\cite{gramfort2013mne} following the pipeline described
in~\cite{poziomska2026quantity}.
A Butterworth high-pass filter (cutoff 0.1\,Hz) removed slow drift,
and a low-pass filter (cutoff 40\,Hz) suppressed high-frequency
muscle artefacts.
A notch filter ($f_0 = 50$\,Hz, quality factor $Q = 5$) attenuated
power-line interference.
All signals were re-referenced to a common average reference and
resampled to 100\,Hz to ensure a uniform sampling rate across all
institutions.
Continuous recordings were segmented into non-overlapping 6-second
frames (600 samples at 100\,Hz).
Frames containing flat-line channels or voltages exceeding
$800\,\mu\text{V}$ were discarded as artefacts.

\subsection{Feature Extraction}
\label{subsec:features}

Three groups of features were extracted from each 6-second frame and
subsequently aggregated per recording using the median across all
valid frames~\cite{poziomska2026quantity}.

\subsubsection{Riemannian Covariance Features}

For each frame, the $19 \times 19$ channel covariance matrix was
computed and averaged using the Riemannian geometric mean via the
\texttt{pyRiemann} library.
The averaged matrix was projected onto its tangent space to obtain a
Euclidean representation; the lower triangular elements were
extracted, yielding \textbf{190 features} per recording.

\subsubsection{Multitaper Power Features}

Power spectral densities were estimated using the multitaper
method~\cite{thomson1982spectrum} and integrated across 14 custom
overlapping frequency bands covering 0.5--40\,Hz:
$[0.5, 2]$, $[1, 3]$, $[2, 4]$, $[3, 6]$, $[4, 8]$, $[6, 10]$,
$[8, 13]$, $[10, 15]$, $[13, 18]$, $[15, 21]$, $[18, 24]$,
$[21, 27]$, $[24, 30]$, and $[27, 40]$\,Hz, spanning the standard
delta, theta, alpha, and beta EEG bands.
PSDs were normalised such that the sum across all channels and bands
equalled one within each frame, yielding \textbf{266 power features}
per recording.

\subsubsection{Inter-Electrode Coherence Features}

Band-wise magnitude-squared coherence was computed for all
$\binom{19}{2} = 171$ inter-electrode pairs in each of the 14
frequency bands:
\begin{equation}
  C_{xy}(f_b) = \frac{\bigl|S_{xy}(f_b)\bigr|^2}
                     {S_{xx}(f_b)\cdot S_{yy}(f_b)},
  \label{eq:coherence}
\end{equation}
where $S_{xy}(f_b)$ is the cross-spectral density and $S_{xx}(f_b)$,
$S_{yy}(f_b)$ are the auto-spectral densities of channels $x$ and $y$
in band $f_b$, yielding $171 \times 14 = \mathbf{2\,394}$
\textbf{coherence features} per recording.

\subsubsection{Final Feature Vector}

Concatenating the three groups yields a final feature vector of
$190 + 266 + 2\,394 = \mathbf{2\,850}$ features per recording,
capturing spatial covariance structure in the time domain, spectral
power across a wide frequency range, and functional connectivity
between all electrode pairs.

\subsection{Approach~1: Feature-Level Harmonisation}
\label{subsec:approach1}

All feature-level methods were applied independently to each of the
2\,850 features.
Harmonisation parameters were estimated exclusively on the training
portion of each data split and subsequently applied to both training
and test data to prevent data leakage.

\subsubsection{Site-Wise Z-Score Standardisation}
\label{subsubsec:sitewise}

As a conceptually simple baseline, each feature was standardised to
zero mean and unit variance using within-site statistics estimated
from the training split:
\begin{equation}
  z_{ij} = \frac{x_{ij} - \hat{\mu}_{s(i),j}}
                {\hat{\sigma}_{s(i),j}},
  \label{eq:sitewise}
\end{equation}
where $s(i)$ is the site at which the recording $i$ was acquired, $\hat{\mu}_{s(i),j}$ and $\hat{\sigma}_{s(i),j}$ are the
within-site mean and standard deviation of feature $j$ estimated on
the training data.
This transformation equalises the first two moments of each feature
distribution across sites but makes no assumptions about the structure
of the site effect, does not account for biological covariates, and
leaves inter-feature covariance structure entirely uncorrected.

\subsubsection{neuroCombat}
\label{subsubsec:neurocombat}

The neuroCombat implementation of the ComBat
algorithm~\cite{johnson2007adjusting,fortin2017harmonization} was
applied with the hospital identifier as the batch variable and patient
age and sex as biological covariates.
ComBat models the observed feature value $y_{ij}$ for subject $i$ and
feature $j$ as
\begin{equation}
  y_{ij} = \alpha_j + \mathbf{x}_i^{\top}\boldsymbol{\beta}_j
           + \gamma_{s(i),j} + \delta_{s(i),j}\,\varepsilon_{ij},
  \label{eq:combat}
\end{equation}
where $\alpha_j$ is the global intercept, $\mathbf{x}_i$ is the
vector of biological covariates with coefficients
$\boldsymbol{\beta}_j$, $\gamma_{s(i),j}$ is an additive site
effect, $\delta_{s(i),j}$ is a multiplicative site effect, and
$\varepsilon_{ij}$ is residual noise, all estimated via empirical
Bayes~\cite{johnson2007adjusting}.
Harmonised values are recovered by:
\begin{equation}
  y^*_{ij} = \frac{y_{ij} - \hat{\alpha}_j
             - \mathbf{x}_i^{\top}\hat{\boldsymbol{\beta}}_j
             - \hat{\gamma}_{s(i),j}}{\hat{\delta}_{s(i),j}}
             + \hat{\alpha}_j
             + \mathbf{x}_i^{\top}\hat{\boldsymbol{\beta}}_j.
  \label{eq:combat_corrected}
\end{equation}


\subsubsection{CovBat}
\label{subsubsec:covbat}

CovBat~\cite{hu2023image} first applies standard ComBat
(Equations~\ref{eq:combat} and~\ref{eq:combat_corrected}) to correct
marginal mean and variance shifts, then performs a second
harmonisation step in the principal-component (PC) space of the
residual covariance matrix.
Specifically, the ComBat-corrected residuals are projected onto their
leading principal components; site-specific shifts in this PC space
are estimated and removed using the ComBat framework; and the
corrected PCs are projected back to the original feature space.
This two-stage procedure targets both marginal and covariance-level
site effects, making CovBat the most comprehensive of the
feature-level methods evaluated here.
We used the implementation of Hu~et~al.~\cite{hu2023image} with
patient age and sex as biological covariates.

\subsubsection{PCA Pre-Orthogonalisation}
\label{subsubsec:pca}

As an additional preprocessing step---applied \emph{before} any of
the marginal correction methods---we evaluated PCA
pre-orthogonalisation of the feature space.
A PCA projection was computed on the training data (retaining 100\%
of variance) and applied to both training and test data prior to
harmonisation.
Because the resulting PC features are mutually uncorrelated by
construction, marginal correction methods subsequently applied to
this transformed space cannot exploit or amplify inter-feature
covariance structure.
PCA pre-orthogonalisation was evaluated both as a standalone
condition and in combination with site-wise standardisation,
neuroCombat, and CovBat.

\subsubsection{Downstream Classifier for Neuroscreening}
\label{subsubsec:classifier}

Original or harmonised feature vectors were passed to a
gradient-boosted (CatBoost~\cite{prokhorenkova2018catboost})
classifier.
Hyperparameters were selected via Optuna~\cite{Akiba2019}
Bayesian optimisation (100 trials, 3-fold stratified CV, AUC
objective) on the unharmonised feature set:
1\,500 boosting iterations, learning rate $0.07$, tree depth $5$,
and $\ell_2$ leaf regularisation $9$.
Early stopping with a patience of 50 rounds was applied during
training.

\subsection{Approach 2: Signal-Level Adversarial Harmonisation}
\label{sec:approach2}

We develop a family of signal-level training conditions based on the
multiple-instance learning architecture of~\cite{poziomska2026quantity},
ranging from a non-adversarial baseline to full Domain Adversarial
Neural Network training~\cite{ganin2016domain} and a subsequent
classifier recalibration variant.
Feature-level harmonisation methods operate on pre-extracted, hand-crafted representations.
Univariate variants such as ComBat correct each feature independently and therefore cannot
address site-related structure residing in inter-feature covariances; CovBat partially
mitigates this by harmonising principal-component scores, yet remains constrained by
linear assumptions and leaves nonlinear inter-feature dependencies unaddressed.
As demonstrated in Section~\ref{subsec:harmonisation_results}, residual site structure
persists in the harmonised representations.
To overcome this, we integrate the adversarial objective with the MIL
architecture of~\cite{poziomska2026quantity}, requiring the network to
solve the pathology classification and site discrimination tasks simultaneously.
Because pathology labels are required for the classification head,
adversarial training is conducted on the full $\mathrm{ELM}_s$ dataset
rather than on the normal-only $\mathrm{ELM}_n$ subset used in the
feature-level experiments.

\subsubsection{Base Architecture}
\label{sec:base_arch}

The base architecture is the MINet model of~\cite{poziomska2026quantity},
consisting of three components (Figure~\ref{fig:dann_arch}).


\begin{figure}[t]
    \centering
    \includegraphics[width=\columnwidth]{Figs/minet_dann_architecture.png}
    \caption{Architecture of the MINet-DANN model.
    The frame encoder~$G_f(\,\cdot\,;\theta_f)$ maps each EEG frame to
$\mathbf{f}_t\!\in\!\mathbb{R}^{288}$; the attention module
($G_a(\,\cdot;\,\theta_a)$) aggregates frame-level representations into
a recording-level descriptor~$\mathbf{h}$.  Two heads branch from
$\mathbf{h}$: the pathology classifier~$G_y$ (top) and the site
discriminator~$G_d$ (bottom).  During back-propagation (dashed
arrows) the Gradient Reversal Layer multiplies the gradient
by~$-\lambda$, encouraging $G_f$~and the attention module to learn
site-invariant features while $G_d$~learns to distinguish sites.}
    \label{fig:dann_arch}
\end{figure}

\paragraph{Frame encoder}
An EEGNet-based~\cite{lawhern2018eegnet} convolutional feature
extractor $G_f(\,\cdot\,;\theta_f)$ maps each 6-second EEG frame
$\mathbf{x}_t \in \mathbb{R}^{C \times T}$ (19~channels, 600~time
points at 100\,Hz) to a 288-dimensional frame-level latent
representation $\mathbf{f}_t \in \mathbb{R}^{288}$:
\begin{equation}
\mathbf{f}_t(\mathbf{x}_t )  = G_f(\mathbf{x}_t;\theta_f).
\label{eq:f_t}
\end{equation}
The encoder comprises 1\,408 parameters and performs separate
temporal and spatial convolutions following~\cite{lawhern2018eegnet}.

\paragraph{Attention-based recording aggregation}
Frame-level representations are aggregated using the non-gated
attention module of~\cite{poziomska2026quantity}, which incorporates
query--key--value structure borrowed from the transformer
architecture.
The module assigns a scalar weight $a_t$ to each frame and computes
the recording-level representation as
\begin{equation}
    \mathbf{h} =  G_a(\,\cdot\,; \theta_a)=\sum_{t=1}^{T_r} a_t\,\mathbf{f}_t, \qquad
    a_t = \frac{%
        \exp\!\bigl(\mathbf{w}^{\top}
        \tanh(\mathbf{V}\mathbf{f}_t)\bigr)}{%
        \sum_{\tau=1}^{T_r}
        \exp\!\bigl(\mathbf{w}^{\top}
        \tanh(\mathbf{V}\mathbf{f}_\tau)\bigr)},
\end{equation}
where $T_r$ is the number of valid frames in the recording and
$\mathbf{V}$, $\mathbf{w}$ are learned parameters.
Frames carrying stronger evidence of pathology are expected to
receive higher attention weights~\cite{poziomska2026quantity}.

\paragraph{Pathology classifier}
The recording-level representation $\mathbf{h}$ is passed to the pathology classifier $G_y(\,\cdot\,;\theta_y)$---a single linear layer acting as the diagnostic head, which produces a binary normal/pathological prediction.

\paragraph{Site discriminator}
A separate two-layer fully connected network $G_d(\,\cdot\,;\theta_d)$, incorporating batch normalization, ReLU activation, and dropout, maps $\mathbf{h}$ to a probability distribution over the 30 recording institutions. Its connection to the backbone differs across the training variants described in Section~\ref{sec:training_variants}: in \mbox{MINet-DANN} it is connected via a Gradient Reversal Layer~(GRL)~\cite{ganin2016domain}, which leaves the input unchanged during forward propagation and multiplies the gradient by $-\lambda$ during back-propagation; in \mbox{MINet-MTL} it is connected via an identity mapping so that gradients flow unmodified; and in the plain MINet baseline, no site discriminator is attached.

\subsubsection{Training Variants}
\label{sec:training_variants}

Four training conditions are evaluated to isolate the contribution
of each component of the adversarial framework.
Their harmonisation effectiveness is assessed on representations
extracted from models trained on the \emph{complete}
$\mathrm{ELM}_s$ dataset
(Section~\ref{subsubsec:two_classifier}), while their
neuroscreening performance is assessed under the
Leave-One-Site-Out (LOSO) protocol
(Section~\ref{subsubsec:clinical_eval}).
This separation ensures that (i)~the harmonisation assessment
reflects the maximum adversarial pressure achievable when all
sites are observed, providing the most direct comparison with the
feature-level site-classification experiments on
$\mathrm{ELM}_n$; and (ii)~the neuroscreening evaluation remains
a rigorous test of generalisation to previously unseen
institutions.

All conditions share the following training protocol.
In each LOSO round, the source pool is partitioned into a training
set and a \emph{source-validation set} (reserving up to 100
recordings per source institution); the entire target site
constitutes the strictly held-out test set.
The first three variants (MINet, MINet-MTL, MINet-DANN) are
initialised from random weights and optimised for a fixed duration
of 50~epochs.
Model selection is performed exclusively on the source-validation
set, preventing any target-domain leakage.
We report performance both at the optimal source-validation epoch
(\textbf{Best}) and at the end of the training cycle
(\textbf{Final}) on the held-out target domain.

\paragraph{MINet (signal-level baseline)}
MINet uses the architecture described above without any domain
objective: the site discriminator $G_d$ is disconnected and
$\lambda \equiv 0$.
It serves as the non-adversarial signal-level reference,
establishing the diagnostic performance achievable by the MIL
backbone in the absence of any harmonisation pressure.
For the separate harmonisation assessment, a single MINet instance
is trained on the complete $\mathrm{ELM}_s$ dataset; the
resulting representations~$\mathbf{h}$ establish the baseline
degree of site information retained without adversarial pressure.

\paragraph{MINet-MTL (multi-task baseline, no gradient reversal)}
MINet-MTL replaces the GRL with an identity mapping, allowing the
gradient from the site discriminator $G_d$ to flow unmodified
into the backbone.
The network is trained under the combined loss:
\begin{equation}
    \mathcal{L}_{\mathrm{MTL}}(\theta_f,\theta_a,\theta_y,\theta_d)
    = 5\cdot\mathcal{L}_y\!\left(G_y(\mathbf{h}),y\right)
      + \mathcal{L}_d\!\left(G_d(\mathbf{h}),d\right),
\end{equation}
where $y$ is the pathology label and $d$ is the site label.
Because the site loss flows normally into the backbone, MINet-MTL
is expected to \emph{encode} rather than suppress site identity,
providing an empirical upper bound on the site information that
an unconstrained multi-task network would retain in~$\mathbf{h}$.

\paragraph{MINet-DANN (full adversarial training)}
MINet-DANN inserts the GRL between the attention module and the
site discriminator, replacing the explicit site-classification
loss with an adversarial objective.
The network is trained under the saddle-point objective:
\begin{equation}
    \mathcal{L}(\theta_f,\theta_a,\theta_y,\theta_d)
    = 5\cdot\mathcal{L}_y\!\left(G_y(\mathbf{h}),y\right)
      - \lambda\,\mathcal{L}_d\!\left(G_d(\mathbf{h}),d\right),
\end{equation}
where the domain-adaptation weight $\lambda$ is annealed from 0
to~1 according to the schedule of~\cite{ganin2016domain}:
\begin{equation}
    \lambda_p = \frac{2}{1+\exp(-\gamma p)} - 1,
    \qquad \gamma = 10,
    \label{eq:lambda_schedule}
\end{equation}
and $p \in [0,1]$ denotes training progress.
This schedule suppresses adversarial pressure in early training,
allowing the encoder and classifier to first develop
discriminative representations before the site-invariance
constraint is fully enforced.
The optimal diagnostic checkpoint is denoted
\emph{MINet-DANN-Best}.

\paragraph{MINet-DANN-CF (classifier fine-tuning on a frozen
adversarial backbone)}
\label{par:MINet-DANN-CF}
MINet-DANN-CF evaluates the diagnostic utility of the
domain-invariant representation learned by MINet-DANN in
isolation from any further representation adaptation.
The network is initialised from the weights of the
\emph{MINet-DANN-Best} checkpoint; the site discriminator $G_d$
is detached, and the entire backbone---comprising the frame
encoder $G_f$ and the attention module (parameters $\theta_f$
and $\theta_a$)---is strictly frozen to preserve the
site-invariant representation~$\mathbf{h}$.
Only the pathology-classifier parameters $\theta_y$ are unfrozen
and optimised using an AdamW optimiser with a reduced learning
rate for 50~epochs, minimising the classification loss alone:
\begin{equation}
    \mathcal{L}(\theta_y)
    = \mathcal{L}_y\!\left(G_y(\mathbf{h}),\,y\right).
\end{equation}
This step recalibrates the diagnostic decision boundaries within
the fixed, adversarially established site-agnostic latent space
without altering the representation itself, directly assessing
whether the adversarially learned~$\mathbf{h}$ is
\emph{intrinsically} sufficient for accurate cross-site pathology
classification.
To mitigate the influence of the original source-side split---which
was used for adversarial-backbone checkpoint selection---the source
pool is randomly re-partitioned into a fresh train/validation
split for the fine-tuning phase, and the best recalibrated state
is selected on this new source-side validation set.

% Four training conditions are evaluated to isolate the contribution of each component of the adversarial framework. Their harmonisation effectiveness is assessed on representations extracted from models trained on the \emph{complete} $\mathrm{ELM}_s$ dataset (Section~\ref{subsubsec:two_classifier}), while their neuroscreening performance is assessed under the Leave-One-Site-Out (LOSO) protocol (Section~\ref{subsubsec:clinical_eval}). This separation ensures that (i)~the harmonisation assessment reflects the maximum adversarial pressure achievable when all sites are observed, providing the most direct comparison with the feature-level site-classification experiments on $\mathrm{ELM}_n$; and (ii)~the neuroscreening evaluation remains a rigorous test of generalisation to previously unseen institutions.

% Crucially, to prevent any target-domain data leakage, all four conditions operate under a \textbf{strict zero-shot LOSO protocol}. In each LOSO round, the available source pool is partitioned into a training set and a \emph{source-validation set} (reserving up to 100 recordings per source institution). Model selection is performed exclusively on this source-validation set. For all variants, we report both the performance at the optimal validation epoch (\textbf{Best}) and at the end of the 50-epoch training cycle (\textbf{Final}) on the strictly unseen target domain.

% \paragraph{MINet (signal-level baseline)}
% The MINet baseline uses the architecture described above without any domain objective (the site discriminator is disconnected, and $\lambda \equiv 0$ in the gradient reversal layer, Figure~\ref{fig:dann_arch}).

% For the neuroscreening evaluation, the network is trained completely from scratch in each LOSO round using exclusively the source pool. The model weights are randomly initialised and optimised using the classification loss for a fixed duration of 50 epochs. Because model selection relies solely on the source-validation subset, the target domain remains entirely unseen at all stages of both representation learning and classifier optimisation. For the separate harmonisation assessment (measuring intrinsic domain clustering), a single instance of MINet is trained on the complete $\mathrm{ELM}_s$ dataset. The resulting recording-level representations $\mathbf{h}$ are extracted to establish the baseline degree of site information retained by the backbone in the absence of any adversarial pressure.

% \paragraph{MINet-MTL (multi-task baseline, no gradient reversal)}
% MINet-MTL replaces the gradient reversal layer (GRL) with an identity mapping, allowing the gradient from the site discriminator $G_d$ to flow unmodified into the backbone. The network is trained from random initialisation using the combined loss:
% \begin{equation}
%     \mathcal{L}_{\mathrm{MTL}}(\theta_f, \theta_a, \theta_y, \theta_d) = 5 \cdot \mathcal{L}_y\!\left(G_y(\mathbf{h}),\,y\right) + \mathcal{L}_d\!\left(G_d(\mathbf{h}),\,d\right),
% \end{equation}
% where $y$ is the pathology label and $d$ is the site label. Because the site loss flows normally into the backbone, MINet-MTL is expected to heavily \emph{encode} site identity rather than suppress it, providing an empirical upper bound on the site information that an unconstrained multi-task network would retain in $\mathbf{h}$. Just like the baseline, it is evaluated under the strict zero-shot LOSO protocol. The absence of adversarial constraints allows the network to freely overfit to source-specific hardware artifacts, which intrinsically explains the high variance typically observed in cross-site performance.

% \paragraph{MINet-DANN (full adversarial training from scratch)}
% MINet-DANN inserts the GRL between the attention module and the site discriminator, replacing the explicit site classification loss with an adversarial objective. The network is trained from random initialisation under the saddle-point objective:
% \begin{equation}
%     \mathcal{L}(\theta_f,\theta_a,\theta_y,\theta_d) = 5 \cdot \mathcal{L}_y\!\left(G_y(\mathbf{h}),\,y\right) - \lambda\,\mathcal{L}_d\!\left(G_d(\mathbf{h}),\,d\right),
% \end{equation}
% where $\theta_f$, $\theta_a$, $\theta_y$, and $\theta_d$ are the parameters of the frame encoder, attention module, pathology classifier, and site discriminator, respectively. The domain-adaptation weight $\lambda$ is annealed from 0 to 1 according to the schedule of \cite{ganin2016domain}:
% \begin{equation}
%     \lambda_p = \frac{2}{1+\exp(-\gamma p)} - 1, \qquad \gamma = 10,
%     \label{eq:lambda_schedule}
% \end{equation}
% where $p \in [0,1]$ denotes training progress. This schedule suppresses adversarial pressure in early training, allowing the encoder and classifier to first develop discriminative representations before the site-invariance constraint is fully enforced. For neuroscreening evaluation, MINet-DANN is trained under the LOSO protocol, and the optimal diagnostic state (MINet-DANN-Best) is captured using the validation AUC, strictly preventing any target-domain leakage.

% \paragraph{MINet-DANN-CF (classifier fine-tuning on a frozen adversarial backbone)}
% \label{par:MINet-DANN-CF}
% MINet-DANN-CF evaluates the diagnostic utility of the domain-invariant representation learned by MINet-DANN in isolation from any further representation adaptation. To achieve this, the backbone is frozen and only the classifier head is
% retrained. We initialize the network using the weights of the optimal \emph{MINet-DANN-Best} model. The site discriminator $G_d$ is detached, and the entire feature extraction backbone---comprising the frame encoder $G_f$ and the attention mechanism (parameters $\theta_f$ and $\theta_a$)---is strictly frozen to preserve the site-invariant representation $\mathbf{h}$. Only the parameters $\theta_y$ of the pathology classifier $G_y$ are unfrozen and optimised.

% Training is performed using an AdamW optimiser with a reduced learning rate for 50 epochs, minimising only the classification loss:
% \begin{equation}
%     \mathcal{L}(\theta_y) = \mathcal{L}_y\!\left(G_y(\mathbf{h}),\,y\right).
% \end{equation}
% This step recalibrates the diagnostic decision boundaries within the fixed, adversarially established site-agnostic latent space without altering the representation itself.

% By decoupling the representation from the classifier, MINet-DANN-CF addresses potential training instabilities observed during prolonged adversarial optimisation, directly assessing whether the adversarially learned $\mathbf{h}$ is \emph{intrinsically} sufficient for accurate cross-site pathology classification. As with all other variants, $G_y$ is fine-tuned exclusively on the source pool. To mitigate the influence of the original source-side split---which was used for adversarial-backbone checkpoint selection---the source pool is
% randomly re-partitioned into a fresh train/validation split for the classifier fine-tuning phase. The best recalibrated state is selected on this new validation set.
% -----------------------------------------------------------------------
\subsubsection{Two-Classifier Site-Discriminability Evaluation}
\label{subsubsec:two_classifier}
The site effect was quantified using a supervised site-classification
task on $\mathrm{ELM}_n$, using 5-fold stratified cross-validation
with stratification by hospital~ID.
To distinguish between the linearly accessible and nonlinearly
accessible components of the site effect, two classifiers of
contrasting capacity were applied to every condition:
\begin{itemize}
  \item \textbf{Logistic regression (LR)} with $\ell_2$
        regularisation, selected by inner cross-validation on the
        training folds.
        LR can only exploit mean shifts and linear inter-feature
        relationships, and therefore measures the component of the
        site effect that survives as a linearly separable signal.
  \item \textbf{CatBoost gradient-boosted classifier}, which can
        additionally exploit higher-order inter-feature interactions
        and covariance structure, and therefore measures the
        \emph{total} site information including structure that is
        invisible to a linear classifier.
\end{itemize}
Performance is reported using the Matthews Correlation Coefficient
(MCC):
\begin{equation}
  \mathrm{MCC} = \frac{c\,s - \sum_{k=1}^{K} p_k t_k}
    {\sqrt{\!\left(s^2 - \sum_{k=1}^{K} p_k^2\right)
           \!\left(s^2 - \sum_{k=1}^{K} t_k^2\right)}},
  \label{eq:mcc}
\end{equation}
where $c = \sum_k C_{kk}$, $s = \sum_{i,j} C_{ij}$,
$t_k = \sum_j C_{kj}$, and $p_k = \sum_i C_{ik}$ for confusion
matrix $\mathbf{C} \in \mathbb{N}^{K \times K}$, and $K$ denotes
the number of classes ($K = 30$ for the site-classification task;
$K = 2$ for the binary neuroscreening task).
The gap between the LR and CatBoost MCC values under any given
condition directly quantifies the degree to which that condition's
site effect is encoded in nonlinear covariance structure.
For feature-level conditions, both classifiers are trained on the
2\,850-dimensional feature vector.
For signal-level conditions, both classifiers are trained on the
recording-level representation $\mathbf{h}$: each signal-level
model is trained once on the complete $\mathrm{ELM}_s$ dataset
(all 30 sites, without held-out folds) and the resulting
representations $\mathbf{h}$ are extracted for all recordings
in $\mathrm{ELM}_n$.
This full-data training condition maximises the adversarial pressure
achievable and therefore represents the strongest achievable
site-invariance of each signal-level approach, providing the most
direct comparison with the feature-level site-classification
experiments.
Note that this harmonisation assessment is kept strictly separate
from the neuroscreening evaluation
(Section~\ref{subsubsec:clinical_eval}), in which signal-level
models are retrained under the LOSO protocol.
Since \mbox{MINet-DANN-CF} does not alter the representation
$\mathbf{h}$ relative to \mbox{MINet-DANN}, both variants share
the same harmonisation assessment result.

\subsubsection{Feature Importance Analysis}
\label{subsubsec:shap}

To identify the primary carriers of site-specific information,
SHAP (SHapley Additive exPlanations) values were computed for the
trained CatBoost site-classification model using the SHAP
TreeExplainer~\cite{lundberg2017unified}.
Global feature importance was summarised by the mean absolute SHAP
value across all predictions and classes.
An ablation study additionally trained models on each feature group
individually (covariance, power, coherence) and with each group
sequentially removed, to corroborate the SHAP-based ranking.

%\subsubsection{Univariate Harmonisation Effectiveness: Cohen's~$d$}
%\label{subsubsec:cohend}

%The impact of harmonisation on individual feature distributions was
%quantified using the standardised mean difference between
%institutions:
%\begin{equation}
%  d = \frac{\mu_1 - \mu_2}{\sigma_{\mathrm{pooled}}},
%  \qquad
%  \sigma_{\mathrm{pooled}} =
%    \sqrt{\frac{(n_1-1)s_1^2 + (n_2-1)s_2^2}{n_1+n_2-2}},
%  \label{eq:cohend}
%\end{equation}
%computed for each feature over every pair of institutions and averaged
%across features to produce an all-vs-all pairwise summary.

\subsubsection{Clinical Neuroscreening Evaluation}
\label{subsubsec:clinical_eval}

Clinical utility was assessed on the binary pathology classification
task using the full $\mathrm{ELM}_s$ dataset under the
Leave-One-Site-Out (LOSO) protocol.
In each round, all recordings from one of the 30 institutions are
held out as the target test set, while the remaining 29 institutions
form the training source pool.
%---
A calibration sample of $n_{\mathrm{cal}} = 30$ normal recordings
from the held-out site is reserved prior to each round.
For feature-level methods, this calibration sample is used together
with the training data from all other institutions to fit the
harmonisation parameters, with the calibration recordings treated as
a separate batch corresponding to the target institution; the
remaining recordings from the target site are then corrected using
the estimated site-specific parameters and constitute the held-out
test set.
For the unharmonised feature-level baseline
the calibration sample is discarded.
For all signal-level conditions (MINet, MINet-MTL, MINet-DANN, and MINet-DANN-CF), the calibration sample is included in the available source pool. Crucially, to prevent any target-domain data leakage during model selection, this source pool is internally partitioned into a training set and a source-validation set. The entire target site's recordings strictly constitute the held-out test set. This protocol directly simulates zero-shot deployment to a previously unseen recording environment and is applied uniformly to all conditions, ensuring direct comparability across the full set of harmonisation strategies.

Performance (MCC and AUC) is reported both as an aggregate across
all 30 rounds and as a per-institution profile for all conditions.

% -----------------------------------------------------------------------
\section{Results}
\label{sec:results}


\subsection{Quantifying the Site Effect}
\label{subsec:site_classification}
The inter-institutional heterogeneity characterised in
Section~\ref{subsec:exploratory}---spanning demographic profiles,
recording durations, amplifier configurations, and raw signal
amplitudes---provides a qualitative basis for expecting strong
site effects in the extracted feature space.
The following sections quantify this effect, evaluate the capacity
of each harmonisation strategy to suppress it, and assess the
downstream impact on clinical neuroscreening performance.

\subsubsection{Unharmonised Site Classification Performance}
The CatBoost nonlinear classifier trained on the unharmonised
$\mathrm{ELM}_n$ feature set achieved an overall MCC of
$0.843 \pm 0.007$ across the 30 institutions in 5-fold stratified
cross-validation, confirming that site identity is strongly and
reliably encoded in the feature space.
The logistic regression linear probe achieved an overall MCC of
$0.785 \pm 0.005$ on the same unharmonised data.
The LR--CatBoost gap of $0.058$ on the unharmonised baseline
establishes the degree to which the site effect in the raw feature
space is already encoded in nonlinear inter-feature dependencies.

Per-institution site-discriminability was further characterised
using one-vs-rest binary MCC for each institution separately
(Table~\ref{tab:per_site_mcc}).
A total of 22 out of 30 institutions achieved one-vs-rest MCC
above 0.7.
A small subset of institutions---H12, H13, H21, and
H23---performed poorly (MCC 0.37--0.48), likely owing to
distributional similarity with neighbouring sites in the corpus.
Institution H5 achieved the lowest one-vs-rest MCC of 0.16.

\begin{table}[htbp]
  \centering
  \caption{Per-institution one-vs-rest MCC (mean $\pm$ SD, 5-fold CV
    on $\mathrm{ELM}_n$, unharmonised CatBoost classifier).
    Institutions are ordered by anonymised identifier.}
  \label{tab:per_site_mcc}
  \begin{tabular}{lc@{\quad}lc}
    \toprule
    \textbf{Institution} & \textbf{MCC} &
    \textbf{Institution} & \textbf{MCC} \\
    \midrule
    H1  & $0.912 \pm 0.017$ & H16 & $0.902 \pm 0.015$ \\
    H2  & $0.743 \pm 0.030$ & H17 & $0.800 \pm 0.033$ \\
    H3  & $0.774 \pm 0.032$ & H18 & $0.869 \pm 0.021$ \\
    H4  & $0.771 \pm 0.042$ & H19 & $0.952 \pm 0.033$ \\
    H5  & $0.160 \pm 0.060$ & H20 & $0.921 \pm 0.030$ \\
    H6  & $0.828 \pm 0.012$ & H21 & $0.451 \pm 0.088$ \\
    H7  & $0.929 \pm 0.016$ & H22 & $0.893 \pm 0.010$ \\
    H8  & $0.805 \pm 0.010$ & H23 & $0.365 \pm 0.081$ \\
    H9  & $0.768 \pm 0.030$ & H24 & $0.764 \pm 0.022$ \\
    H10 & $0.872 \pm 0.006$ & H25 & $0.887 \pm 0.017$ \\
    H11 & $0.864 \pm 0.061$ & H26 & $0.785 \pm 0.034$ \\
    H12 & $0.480 \pm 0.096$ & H27 & $0.715 \pm 0.020$ \\
    H13 & $0.369 \pm 0.142$ & H28 & $0.924 \pm 0.031$ \\
    H14 & $0.939 \pm 0.012$ & H29 & $0.663 \pm 0.014$ \\
    H15 & $0.856 \pm 0.009$ & H30 & $0.547 \pm 0.098$ \\
    \midrule
    \multicolumn{2}{l}{overall} &
    \multicolumn{2}{l}{$0.843 \pm 0.007$} \\
    \bottomrule
  \end{tabular}
\end{table}

\subsubsection{Feature Importance: SHAP Analysis and Ablation}
\label{subsubsec:shap_results}

Table~\ref{tab:ablation} summarises site-classification MCC for
CatBoost models trained on individual feature groups and with each
group excluded.
A model trained exclusively on coherence features achieved MCC of
$0.850 \pm 0.003$, nearly matching the full-feature-set performance.
Power features alone achieved $0.600 \pm 0.005$, and covariance
features alone achieved $0.752 \pm 0.009$.
Excluding coherence features reduced performance most substantially
(to $0.785 \pm 0.007$), while excluding power or covariance features
had comparatively little effect, confirming that coherence features
are the dominant carriers of site-specific information.

The SHAP analysis corroborates these findings.
Among the coherence features, the O1--O2 electrode pair achieved the
highest importance scores across the majority of frequency bands.
Beta-band coherence was particularly prominent in the temporal regions
(T3, T4, T5, T6, F7, F8).
Among spectral power features, delta-band power and the power at
electrode Fz were the most important.
Covariance features contributed comparatively little to the model's
site predictions.

\begin{table}[htbp]
  \centering
  \caption{Overall site-classification MCC (mean $\pm$ SD, 5-fold CV
    on $\mathrm{ELM}_n$, CatBoost classifier) for the feature-group
    ablation study.}
  \label{tab:ablation}
  \begin{tabular}{lcc}
    \toprule
    \textbf{Condition} & \textbf{MCC} & \textbf{SD} \\
    \midrule
    \multicolumn{3}{l}{\textit{Single feature group}} \\
    \quad Coherence only  & 0.850 & 0.003 \\
    \quad Power only      & 0.600 & 0.005 \\
    \quad Covariance only & 0.752 & 0.009 \\
    \midrule
    \multicolumn{3}{l}{\textit{One group excluded}} \\
    \quad All except coherence  & 0.785 & 0.007 \\
    \quad All except power      & 0.857 & 0.002 \\
    \quad All except covariance & 0.863 & 0.004 \\
    \bottomrule
  \end{tabular}
\end{table}


\subsection{The Two-Classifier Evaluation of Harmonisation Methods}
\label{subsec:harmonisation_results}

As a sanity check on marginal correction, mean pairwise Cohen's~$d$
(averaged across all institution pairs and all 2\,850 features)
was reduced from [TBF] in the unharmonised baseline to [TBF]
after site-wise standardisation, [TBF] after \mbox{neuroCombat},
and [TBF] after \mbox{CovBat}, confirming that all three methods
successfully equalise marginal feature distributions across sites.
The two-classifier evaluation below reveals, however, that this
univariate reduction is insufficient as a sole indicator of
harmonisation quality.


\subsubsection{The Harmonisation Paradox Revealed by the
  Two-Classifier Framework}
\label{subsubsec:paradox}

Table~\ref{tab:site_class_harmonisation} reports the LR-MCC and
CatBoost-MCC for all conditions.
The logistic regression linear probe assigned near-chance MCC to all
harmonised conditions, confirming that marginal correction
successfully removes the \emph{linear} component of the site effect. However, CovBat retained a small linear residual (LR-MCC $= 0.176$) - this could suggest that its covariance-level correction leaves some linear site signal
The CatBoost nonlinear probe, by contrast, assigned near-ceiling MCC
($0.924$--$0.969$) to all harmonised conditions applied without PCA
pre-orthogonalisation, demonstrating the \emph{harmonisation paradox}:
the nonlinear covariance component of the site effect not only
survives marginal correction but becomes \emph{more} salient to the
nonlinear classifier.

The LR--CatBoost gap is therefore the quantitative fingerprint of the harmonisation paradox: a large gap indicates that the residual site effect after correction is encoded predominantly in nonlinear inter-feature dependencies rather than in marginal feature distributions.
PCA pre-orthogonalisation substantially reduces this gap by
destroying the inter-feature covariance prior to correction,
attenuating the paradox across all harmonisation conditions.

\td{
Grouped bar chart z dwiema grupami słupków (LR i CB) per warunek, posortowanymi w kolejności rosnącego gapu:
Oś X: warunki harmonizacji pogrupowane w trzy bloki wizualnie oddzielone odstępem lub kolorem tła:
Blok 1: warunki bez PCA (Unharmonised, Site-wise, neuroCombat, CovBat)
Blok 2: warunki z PCA (PCA only, PCA+Site-wise, PCA+neuroCombat, PCA+CovBat)
Blok 3: signal-level (MINet, MINet-MTL, MINet-DANN)
Dla każdego warunku harmonizacji rysowane są dwa sąsiadujące słupki — jeden dla LR-MCC i jeden dla CB-MCC — z whiskerami SD. Warunki są zgrupowane w trzy bloki jak opisano wyżej, oddzielone wizualnie węższymi przerwami między blokami niż między parami słupków wewnątrz bloku.
Między każdą parą słupków LR i CB można dodać cienkę dwukierunkową strzałkę lub klamrę z wartością ΔMCC
ΔMCC wpisaną bezpośrednio, żeby gap był czytelny numerycznie bez potrzeby patrzenia do tabeli. Alternatywnie wartości Δ można umieścić jako osobną linię punktów nałożoną na wykres słupkowy (secondary axis lub po prostu punkty nad parami słupków).
}

The signal-level conditions are evaluated via the same two-classifier
framework applied to the representation $\mathbf{h}$.
\textbf{[TBF: report LR-MCC and CatBoost-MCC for all four
signal-level variants once experiments are complete, and comment on
whether MINet-DANN achieves a smaller LR--CatBoost gap than the
feature-level conditions.]}

Figure~\ref{fig:site_mcc_no_pca} shows per-institution
site-classification MCC (CatBoost) without PCA
pre-orthogonalisation, and Figure~\ref{fig:site_mcc_pca} shows the
analogous results with PCA pre-orthogonalisation.

\begin{figure}[htbp]
  \centering
  \includegraphics[width=\textwidth]{Figs/site_mcc_pca_00.pdf}
  \caption{Per-institution site-classification MCC (CatBoost,
    one-vs-rest, 5-fold CV on $\mathrm{ELM}_n$) for the raw feature
    set and after four harmonisation methods \emph{without} PCA
    pre-orthogonalisation.
    Institutions are ordered by decreasing raw-baseline MCC.
    All four methods drive MCC towards a near-ceiling value,
    demonstrating the harmonisation paradox.}
  \label{fig:site_mcc_no_pca}
\end{figure}

\begin{figure}[htbp]
  \centering
  \includegraphics[width=\textwidth]{Figs/site_mcc_pca_10.pdf}
  \caption{Per-institution site-classification MCC (CatBoost,
    one-vs-rest, 5-fold CV on $\mathrm{ELM}_n$) after four
    harmonisation methods \emph{with} full-variance PCA
    pre-orthogonalisation.
    Institutions are ordered by decreasing raw-baseline MCC.
    PCA pre-orthogonalisation substantially attenuates the
    harmonisation paradox.}
  \label{fig:site_mcc_pca}
\end{figure}

\begin{table}[htbp]
  \centering
  \caption{Overall site-classification MCC (mean $\pm$ SD, 5-fold CV
    on $\mathrm{ELM}_n$) for all harmonisation conditions under the
    two-classifier framework.
    LR = logistic regression linear probe; CB = CatBoost nonlinear
    probe.
    The LR--CB gap quantifies the degree to which residual site
    information is encoded in nonlinear covariance structure.
    $^\dagger$Signal-level models: site discriminability assessed on
    representations from full-data training
    (Section~\ref{subsubsec:repr_assessment}).
    MINet-DANN-CF shares the harmonisation assessment of MINet-DANN
    as classifier retraining does not alter~$\mathbf{h}$.}
  \label{tab:site_class_harmonisation}
  \begin{tabular}{lcccc}
    \toprule
    & \multicolumn{2}{c}{\textbf{LR (linear)}}
    & \multicolumn{2}{c}{\textbf{CatBoost (nonlinear)}} \\
    \cmidrule(lr){2-3}\cmidrule(lr){4-5}
    \textbf{Condition} & MCC & SD & MCC & SD \\
    \midrule
    \multicolumn{5}{l}{\textit{Feature-level, without PCA}} \\
    \quad Original (unharmonised)
     & 0.785 & 0.005 & 0.843 & 0.007 \\
    \quad Site-wise standardisation
      & 0.000 & 0.000 & 0.969 & 0.002 \\
    \quad neuroCombat
      & 0.013 & 0.006 & 0.966 & 0.002 \\
    \quad CovBat
      & 0.176 & 0.013 & 0.924 & 0.004 \\
    \midrule
    \multicolumn{5}{l}{\textit{Feature-level, with PCA
      pre-orthogonalisation}} \\
    \quad PCA only
      & 0.858 & 0.004 & 0.784 & 0.008 \\
    \quad PCA + site-wise stand.
      & 0.000 & 0.000 & 0.720 & 0.021 \\
    \quad PCA + neuroCombat
      & 0.034 & 0.005 & 0.710 & 0.016 \\
    \quad PCA + CovBat
      & 0.032 & 0.007 & 0.740 & 0.014 \\
    \midrule
    \multicolumn{5}{l}{\textit{Signal-level (both probes on
      $\mathbf{h}$)$^\dagger$}} \\
    \quad MINet ($\lambda \equiv 0$)
      & \textbf{[TBF]} & \textbf{[TBF]}
      & \textbf{[TBF]} & \textbf{[TBF]} \\
    \quad MINet-MTL
      & \textbf{[TBF]} & \textbf{[TBF]}
      & \textbf{[TBF]} & \textbf{[TBF]} \\
    \quad MINet-DANN
      & \textbf{[TBF]} & \textbf{[TBF]}
      & \textbf{[TBF]} & \textbf{[TBF]} \\
      % \quad MINet-DANN-CF
      % & \textbf{[TBF]} & \textbf{[TBF]}
      % & \textbf{[TBF]} & \textbf{[TBF]} \\
    \bottomrule
  \end{tabular}
\end{table}




% -----------------------------------------------------------------------
\subsection{Impact of Harmonisation on Clinical Neuroscreening}
\label{subsec:clinical_results}

\subsubsection{Per-Institution LOSO Performance}

Figure~\ref{fig:patho_auc_no_pca} shows the per-institution
neuroscreening AUC under the LOSO protocol for all feature-level
harmonisation methods applied without PCA pre-orthogonalisation.
Without PCA, the four methods diverge noticeably at individual
institutions.
\textbf{[TBF: identify specific outlier institutions, describe the
direction of the anomaly, and comment on likely causes.]}

Figure~\ref{fig:patho_auc_pca} shows the analogous results with PCA
pre-orthogonalisation.
Under this condition, the four methods cluster closely together across
all institutions, and the institution-level variability in AUC is
markedly reduced relative to the no-PCA setting.

\begin{figure}[htbp]
  \centering
  \includegraphics[width=\textwidth]{Figs/patho_auc_pca_none.pdf}
  \caption{Per-institution neuroscreening AUC under the LOSO protocol
    for the CatBoost classifier applied to raw features and after four
    feature-level harmonisation methods, \emph{without} PCA
    pre-orthogonalisation.
    Institutions are ordered by decreasing raw-baseline AUC.}
  \label{fig:patho_auc_no_pca}
\end{figure}

\begin{figure}[htbp]
  \centering
  \includegraphics[width=\textwidth]{Figs/patho_auc_pca_all.pdf}
  \caption{Per-institution neuroscreening AUC under the LOSO protocol
    for the CatBoost classifier applied to raw features and after four
    feature-level harmonisation methods, \emph{with} full-variance PCA
    pre-orthogonalisation.
    Institutions are ordered by decreasing raw-baseline AUC.
    All four methods cluster closely and track slightly above the raw
    baseline, with substantially reduced institution-level variability
    compared to the no-PCA condition
    (Figure~\ref{fig:patho_auc_no_pca}).}
  \label{fig:patho_auc_pca}
\end{figure}

\subsubsection{Signal-Level Conditions}
\label{subsubsec:dann_clinical}
Before reporting neuroscreening performance, we verify that the
adversarial objective achieved its intended effect on the
recording-level representation~$\mathbf{h}$.
MINet-MTL, which trains the site discriminator without gradient
reversal, retained strong site-discriminative information in~$\mathbf{h}$ (domain MCC $= 0.568 \pm 0.077$), confirming that an unconstrained
multi-task objective actively encodes site identity.
MINet-DANN, by contrast, drove domain MCC to near-zero
($0.009 \pm 0.014$, Figure~\ref{fig:dann_domain_mcc}),
confirming that the Gradient Reversal Layer successfully rendered
the latent space site-agnostic.
Because MINet-DANN-CF operates on the same frozen backbone, its
domain MCC is identical by design.

Under the strict zero-shot protocol, the optimal MINet-DANN model (selected via the source-validation set) achieved a mean target AUC of $0.791 \pm 0.052$ and a mean MCC of $0.452 \pm 0.105$ across all 30 LOSO institutions.

However, MINet-DANN also exhibited a training stability problem: target-domain performance degraded at the end of the training cycle, with the final-epoch validation AUC falling to $0.777 \pm 0.056$, visibly below the best-epoch AUC (Figure~\ref{fig:stability_scatter_dann}). The observed decline in target-domain AUC beyond the optimal epoch is consistent with two competing explanations: catastrophic forgetting, in which the adversarial objective progressively disrupts diagnostically relevant structure in the latent representation $\mathbf{h}$, and standard overfitting, in which the pathology classifier memorizes source-pool-specific patterns that do not transfer to the target institution. These can be distinguished by monitoring source-pool validation AUC alongside target AUC: a simultaneous decline in both metrics implicates representation-level disruption, whereas a divergence—source AUC stable or improving while target AUC falls—indicates overfitting. Figure~\ref{fig:stability_scatter_dann} plots target AUC only; disentangling these mechanisms is left for future work.

MINet-DANN-CF directly addresses this instability by freezing the adversarially trained backbone at its best checkpoint and retraining only the pathology classifier $G_y$ on the source pool under the same strict LOSO protocol. This strategy successfully recovered and stabilized diagnostic performance, achieving a highly robust final-epoch mean AUC of $0.834 \pm 0.041$ and a mean MCC of $0.498 \pm 0.106$. This marks a substantial quantitative improvement over the raw MINet-DANN and even outperforms the unharmonized Baseline (Best AUC $0.822 \pm 0.043$), confirming that the adversarially scrubbed latent representations are inherently highly discriminative when decoupled from the instability of joint adversarial training.

Per-institution AUC profiles for MINet, MINet-DANN, and MINet-DANN-CF are shown in Figure~\ref{fig:all_strategies_auc} ( \ref{fig:all_strategies_mcc} shows analogous profiles for MCC).


\begin{figure}[htbp]
  \centering
  \includegraphics[width=\textwidth]{manuscript_Biomedical_Signal_Processing_and_Control/Figs/auc_new_labels_h1-30.png}
\caption{Per-institution neuroscreening performance (AUC) under the Leave-One-Site-Out (LOSO) protocol. The figure compares four signal-level training conditions: MINet (unharmonised baseline), MINet-MTL, MINet-DANN, and MINet-DANN-CF. Performance is reported for the optimal model checkpoint (Best epoch). Crucially, to prevent any target-domain data leakage, this optimal state was selected exclusively via a held-out source-validation set. This ensures that the reported results represent a strict zero-shot assessment on the previously unseen target institution.}
  \label{fig:all_strategies_auc}
\end{figure}



\begin{figure}[htbp]
  \centering
  \includegraphics[width=0.6\textwidth]{manuscript_Biomedical_Signal_Processing_and_Control/Figs/dann_training_stability_new_labels.png}
  \caption{Training stability for MINet-DANN: best-epoch vs.\
    final-epoch target AUC across all 30 LOSO rounds.
    Points below the diagonal indicate performance degradation
     at the end of the training cycle.}
  \label{fig:stability_scatter_dann}
\end{figure}

\newpage
\subsubsection{Aggregate Neuroscreening Performance}

Table~\ref{tab:clinical_loso} consolidates aggregate neuroscreening
performance across all feature-level and signal-level conditions
under the LOSO protocol.
Known results (unharmonised, site-wise standardisation, and
neuroCombat) are included where available; all other cells await
completion of the extended experiments.
For reference, preliminary 5-fold cross-validation results on these
three baseline conditions confirmed that harmonisation benefits are
not confined to the zero-shot setting: site-wise standardisation and
neuroCombat each improved CV AUC by approximately 0.015 relative to
the unharmonised baseline, consistent with the direction and
magnitude of the LOSO gains reported here.

\textbf{[TBF: narrative summary once all conditions are complete,
addressing: (i) whether harmonisation consistently improves on the
unharmonised baseline; (ii) additional aggregate benefit of PCA
pre-orthogonalisation; (iii) relative performance of CovBat versus
neuroCombat and site-wise standardisation; (iv) where the
signal-level conditions sit relative to feature-level conditions.]}

\begin{table}[htbp]
  \centering
  \caption{Neuroscreening performance under the Leave-One-Site-Out
    protocol (mean $\pm$ SD across 30 rounds) for all conditions.
    }
  \label{tab:clinical_loso}
  \begin{tabular}{lcc}
    \toprule
    \textbf{Method} & \textbf{MCC} & \textbf{AUC} \\
    \midrule
    \multicolumn{3}{l}{\textit{CatBoost, without PCA
      pre-orthogonalisation}} \\
    \quad Unharmonised
      & $0.399 \pm 0.134$ & $0.800 \pm 0.078$ \\
    \quad Site-wise standardisation
      & $0.407 \pm 0.116$ & $0.801 \pm 0.064$ \\
    \quad neuroCombat
      & $0.420 \pm 0.122$ & $0.808 \pm 0.072$ \\
    \quad CovBat
      & $0.392 \pm 0.128$ & $0.791 \pm 0.070$ \\
    \midrule
    \multicolumn{3}{l}{\textit{CatBoost, with PCA
      pre-orthogonalisation}} \\
    \quad PCA only
      & $0.370 \pm 0.124$ & $0.782 \pm 0.076$ \\
    \quad PCA + site-wise stand.
      & $0.389 \pm 0.127$ & $0.788 \pm 0.080$ \\
    \quad PCA + neuroCombat
      & $0.391 \pm 0.129$ & $0.787 \pm 0.079$ \\
    \quad PCA + CovBat
      & $0.395 \pm 0.127$ & $0.789 \pm 0.080$ \\
    \midrule
    \multicolumn{3}{l}{\textit{Logistic Regression, without PCA pre-orthogonalisation}} \\
    \quad Unharmonised
      & $0.359 \pm 0.124$ & $0.770 \pm 0.079$ \\
    \quad Site-wise standardisation
      & $0.373 \pm 0.115$ & $0.769 \pm 0.066$ \\
    \quad neuroCombat
      & $0.382 \pm 0.122$ & $0.775 \pm 0.076$ \\
    \quad CovBat
      & $0.386 \pm 0.123$ & $0.783 \pm 0.070$ \\
    \midrule
    \multicolumn{3}{l}{\textit{Logistic Regression, with PCA pre-orthogonalisation}} \\
    \quad PCA only
      & $0.356 \pm 0.126$ & $0.771 \pm 0.078$ \\
    \quad PCA + site-wise stand.
      & $0.391 \pm 0.118$ & $0.782 \pm 0.075$ \\
    \quad PCA + neuroCombat
      & $0.383 \pm 0.115$ & $0.782 \pm 0.075$ \\
    \quad PCA + CovBat
      & $0.386 \pm 0.119$ & $0.783 \pm 0.076$ \\
    \midrule
    \multicolumn{3}{l}{\textit{Signal-level conditions}} \\
    \quad MINet ($\lambda \equiv 0$)
      & $0.501 \pm 0.080$ & $0.822 \pm 0.043$ \\
    \quad MINet-MTL
      & $0.400 \pm 0.096$ & $0.763 \pm 0.058$ \\
    \quad MINet-DANN
      & $0.452 \pm 0.105$ & $0.791 \pm 0.052$ \\
    \quad MINet-DANN-CF
      & $0.498 \pm 0.106$ & $0.834 \pm 0.041$ \\
    \bottomrule
  \end{tabular}
\end{table}





% -----------------------------------------------------------------------
\section{Discussion}
\label{sec:discussion}
\td{dyskusja do dokładniejszego sczytania jak już będzie komplet wyników}
\subsection{The Nature and Magnitude of the Site Effect}
\label{subsec:disc_site_effect}

The site-classification experiment confirmed that the site effect is
both strong and broadly distributed across the 2\,850-dimensional
feature space of $\mathrm{ELM}_n$, with the CatBoost nonlinear
classifier achieving an overall MCC of $0.843 \pm 0.007$.
The range of per-site MCC values---from 0.952 (H19) to 0.160
(H5)---reflects the heterogeneity of the corpus: institutions with
distinctive demographic profiles (e.g.\ H7 and H19, which serve
predominantly older patients) or unusual recording protocols
(e.g.\ H16, with sessions up to 8\,322 seconds) are most reliably
identified.
Institutions at the lower end of the MCC range (H5, H13, H21, H23)
are likely to have feature distributions that overlap substantially
with those of neighbouring sites, either because of similar patient
demographics or because their recording protocols are less
distinctive.

The SHAP analysis and ablation study converge on a consistent
interpretation: inter-electrode coherence features are the primary
carriers of site-specific information.
The elevated importance of O1--O2 coherence is consistent with the
hypothesis that differences in eye-open versus eye-closed protocols
across institutions modulate occipital alpha coherence, which is
highly sensitive to visual state.
The prominence of beta-band coherence in temporal regions is
consistent with variability in muscle artefact levels arising from
differences in patient preparation, scalp cleaning, and electrode
impedance management across sites, since the beta band overlaps with
the electromyographic spectrum.
Delta-band power and frontal Fz power, identified as important
spectral components, may reflect differences in patient alertness and
vigilance state at recording time.

\subsection{The Harmonisation Paradox and the Two-Classifier
  Interpretation}
\label{subsec:disc_paradox}

The most striking finding of this study is the \emph{harmonisation
paradox} (Section~\ref{subsubsec:paradox}): all marginal correction
methods successfully reduced pairwise Cohen's~$d$ inter-site
differences and drove the logistic regression linear probe to
near-chance site-classification MCC, yet simultaneously increased the
CatBoost nonlinear probe MCC towards a near-ceiling value.

The two-classifier framework resolves this apparent contradiction.
Marginal correction adjusts the mean and variance of each feature
independently, thereby equalising marginal distributions across sites
and removing the linear component of the site effect.
However, the site effect is encoded not only in marginal distributions
but also in the \emph{covariance structure} between features---a form
of inter-site variability that feature-by-feature corrections leave
entirely intact~\cite{hu2023image}.
By equalising means and variances, these methods remove the dominant
sources of within-site variance, causing the subtler but highly
structured inter-feature covariance patterns that encode site identity
to become more salient to a nonlinear classifier.
This argument is directly analogous to that advanced by
Mellot~et~al.~\cite{mellot2023meeg} in the context of Riemannian
geometry-based EEG alignment.

The PCA pre-orthogonalisation experiments provide direct mechanistic
evidence: by projecting the feature space onto mutually uncorrelated
principal components, PCA eliminates inter-feature covariance in the
original feature basis.
Marginal correction applied to this PC-transformed space therefore
cannot amplify covariances in the original basis, and the gap between
LR-MCC and CatBoost-MCC narrows substantially.

This finding has a direct methodological implication: univariate
metrics such as Cohen's~$d$ and the linear probe MCC are inadequate
as sole indicators of harmonisation quality in high-dimensional
feature spaces.
A complete evaluation framework must include a nonlinear
site-discriminability measure---such as the CatBoost MCC reported
here---as a necessary complement to feature-level effect-size
summaries.

\subsection{Signal-Level Adversarial Harmonisation}
\label{subsec:disc_dann}
The four signal-level training conditions evaluated in this study
provide converging evidence on both the effectiveness and the
limitations of end-to-end adversarial harmonisation in the clinical
EEG setting.

The domain MCC results confirm that the Gradient Reversal Layer
successfully decouples site-specific signatures from the
recording-level representation~$\mathbf{h}$.
MINet-MTL, which trains the site discriminator without gradient
reversal, actively encodes site identity---as expected for an
unconstrained multi-task objective---and serves as the empirical
upper bound on the site information that such a network would retain.
MINet-DANN achieves near-zero domain MCC on the source pool,
confirming that the adversarial objective forced the latent space to
become site-agnostic.

A key practical finding concerns the training instability of
MINet-DANN: the adversarial objective gradually distorts the
diagnostic decision boundaries, causing final-epoch performance to
fall substantially below the best-epoch AUC in many LOSO rounds.
Whether this reflects true catastrophic forgetting---in which the
adversarial objective progressively disrupts diagnostically relevant
structure in~$\mathbf{h}$---or overfitting of the diagnostic head
to source-pool-specific patterns that do not transfer to the target
institution cannot be determined from target AUC alone and warrants
further investigation.
MINet-DANN-CF directly addresses this instability by decoupling the
representation from the classifier: the adversarially trained backbone
is frozen at its best checkpoint, and only the pathology classifier
$G_y$ is retrained on the source pool under the same strict zero-shot
LOSO protocol.
Because the representation $\mathbf{h}$ is unchanged, the
site-invariance established by MINet-DANN is fully preserved, while
the classifier head is recalibrated to the site-agnostic feature
space without further adversarial distortion.

The key empirical question is whether MINet-DANN-CF recovers the
diagnostic performance lost during adversarial training and whether
it improves over both the MINet baseline and the feature-level LOSO
conditions.
\textbf{[TBF: complete this comparison once MINet-DANN-CF LOSO
results are available, and address whether the LR--CatBoost gap on
$\mathbf{h}$ for MINet-DANN is smaller than for all feature-level
conditions, indicating suppression of the nonlinear covariance
component of the site effect.]}
\subsection{Impact on Clinical Neuroscreening}
\label{subsec:disc_clinical}


Under the strict zero-shot LOSO protocol, the optimal MINet-DANN model achieved a mean diagnostic AUC of $0.791 \pm 0.052$ and an MCC of $0.452 \pm 0.105$. While this performance surpasses the naive multi-task approach (MINet-MTL, Best AUC $0.763 \pm 0.058$), it falls short of the unharmonized signal-level baseline (Best AUC $0.822 \pm 0.043$). However, when the adversarially harmonized representations are decoupled from the joint training process via classifier finetuning, MINet-DANN-CF recovers and exceeds this performance, achieving the highest overall diagnostic accuracy with a robust best-epoch AUC of $0.834 \pm 0.041$ and an MCC of $0.498 \pm 0.106$ highly comparable to final-epoch AUC of $0.835 \pm 0.042$ and an MCC of $0.499 \pm 0.104$.
\textbf{[TBF: complete once CovBat, PCA, MINet, MINet-MTL, and
MINet-DANN-CF LOSO results are available.]}

\subsection{Limitations}
\label{subsec:limitations}

Several limitations of the present study should be noted.
First, although patient age and sex were included as biological
covariates in neuroCombat and CovBat, other clinically relevant
variables---such as the reason for referral, medication status, or
cognitive state at recording time---were not available for the full
corpus and could not be controlled for.
Second, all marginal correction methods adjust only the first two
moments of each feature's distribution and cannot correct
higher-order distributional differences.
Third, PCA pre-orthogonalisation discards the covariance structure
of the original feature space; while this attenuates the
harmonisation paradox, it also removes inter-feature relationships
that may carry clinically relevant information, and the long-term
impact on downstream task performance warrants further investigation.
Fourth, the CatBoost probe applied to the DANN representations
$\mathbf{h}$ evaluates nonlinear site discriminability in a
288-dimensional latent space rather than the original
2\,850-dimensional feature space; while this is internally consistent
within the two-classifier framework, direct numerical comparison
between the feature-level and signal-level CatBoost-MCC values should
be interpreted with caution.
Fifth, although MINet-DANN-CF eliminates the catastrophic forgetting
observed in MINet-DANN by freezing the backbone, it depends on the
quality of the MINet-DANN checkpoint selected; if the best-epoch
adversarial checkpoint does not yet achieve sufficient site-invariance,
the recalibrated classifier inherits a partially site-confounded
representation.
Sixth, the present evaluation is confined to the LOSO protocol, which
measures generalisation to completely unseen institutions.
The performance of the evaluated harmonisation methods in the
interpolation setting---where all institutions are represented in
training---may differ, and this complementary perspective warrants
investigation in future work.

\subsection{Future Directions}
\label{subsec:future}

The harmonisation paradox and the two-classifier framework identified
in this study motivate several directions for future work.
Most immediately, the completed signal-level results will determine
whether adversarial training can suppress nonlinear covariance-based
site structure, which neither marginal correction nor PCA
pre-orthogonalisation can fully address.
Beyond harmonisation methodology, the SHAP analysis suggests that
feature engineering strategies that reduce coherence-feature
sensitivity to site-specific procedural differences could reduce the
effective site effect without explicit harmonisation.
Finally, normative modelling frameworks~\cite{dong2025normative},
which express individual measurements as deviations from an
age-adjusted population norm, constitute a complementary approach
that may interact favourably with the harmonisation methods studied
here and warrant systematic evaluation in the $\mathrm{ELM}_s$
context.

% -----------------------------------------------------------------------
\section{Conclusion}
\label{sec:conclusion}

This study provides a comprehensive characterisation of the site
effect in a large-scale clinical EEG corpus ($\mathrm{ELM}_s$,
54\,779 recordings, 30 hospitals) and a systematic evaluation of
feature-level and signal-level harmonisation strategies under a
unified two-classifier evaluation framework.
All conditions are evaluated under a rigorous Leave-One-Site-Out
protocol, directly simulating zero-shot deployment to previously
unseen recording environments.

The CatBoost nonlinear classifier achieved a site-identification MCC
of $0.843 \pm 0.007$, confirming that site identity is strongly
encoded in EEG feature representations.
SHAP explainability analysis and ablation experiments identified
inter-electrode coherence---particularly O1--O2 and beta-band
temporal coherence---as the dominant carrier of site-specific
information, implicating differences in eye-state protocols and
muscle artefact management across institutions.

A central finding is the \emph{harmonisation paradox}: all marginal
correction methods successfully removed the linear component of the
site effect (Cohen's~$d$ reduced; logistic regression MCC near
chance) yet simultaneously amplified the nonlinear CatBoost MCC to
near-ceiling values, demonstrating that feature-by-feature correction
reorganises rather than eliminates the multivariate covariance
structure of the site effect.
PCA pre-orthogonalisation attenuates this paradox and yields more
uniform and stable cross-site neuroscreening performance.
These findings demonstrate that a complete harmonisation evaluation
framework must include a nonlinear site-discriminability measure
alongside univariate metrics.

The necessity of the adversarial framework is starkly highlighted when comparing MINet-DANN against the naive multi-task approach (MINet-MTL). Without the gradient reversal mechanism to actively penalize site-specific features, MINet-MTL explicitly learns to identify the recording institutions, reaching a high Domain MCC of $0.544 \pm 0.079$ by the end of training. This creates a catastrophic feature trade-off: as the network increasingly encodes hardware and protocol artifacts into the latent space, it actively sacrifices diagnostically relevant information. Consequently, its zero-shot diagnostic performance collapses, with the final-epoch target AUC plummeting to $0.666 \pm 0.063$ (down from an earlier, already sub-optimal validation peak of $0.763 \pm 0.058$).
In contrast, signal-level adversarial harmonisation via MINet-DANN successfully eliminated this site-specific bias (reducing Final Domain MCC to near zero: $0.006 \pm 0.006$) while restoring diagnostic efficacy across the 30 institutions. Under the strict zero-shot LOSO protocol, the optimal MINet-DANN model achieved a mean diagnostic AUC of $0.791 \pm 0.052$ and an MCC of $0.452 \pm 0.105$. While this performance surpasses the naive multi-task approach (MINet-MTL, Best AUC $0.763 \pm 0.058$), it falls short of the unharmonized signal-level baseline (Best AUC $0.822 \pm 0.043$). However, when the adversarially harmonized representations are decoupled from the joint training process, MINet-DANN-CF recovers and exceeds this performance. It achieves the highest overall diagnostic accuracy while demonstrating exceptional training stability, with the optimal validation state (Best AUC $0.834 \pm 0.041$, MCC $0.498 \pm 0.106$) being virtually indistinguishable from the highly robust end-of-training state (Final AUC $0.835 \pm 0.042$, MCC $0.499 \pm 0.104$).

Building upon this, MINet-DANN-CF aims to optimally recalibrate the diagnostic decision boundaries using only the source pool, ensuring the site-invariant representation remains strictly preserved. Freezing the adversarially trained backbone and retraining solely the pathology classifier yielded our highest aggregate performance under the strict zero-shot protocol, achieving a robust final-epoch mean AUC of $0.835 \pm 0.042$ and an MCC of $0.499 \pm 0.104$. This confirms the high diagnostic utility of the domain-invariant features extracted by the adversarial framework, proving they remain highly discriminative and stable even when decoupled from the joint adversarial training process.

These results collectively underscore that the site effect is a
complex, high-dimensional challenge for which standard marginal
correction methods provide only a partial solution.
Developing truly generalisable neuroscreening models will require
harmonisation strategies that explicitly account for multivariate
inter-site structure---whether through covariance-aware statistical
corrections, adversarial representation learning, or feature
engineering guided by explainability analysis.

% -----------------------------------------------------------------------
\section*{Acknowledgements}
This study was supported by The National Centre for Research and
Development (NCBR) grant number POIR.01.01.01-00-1074/21

\section*{CRediT Author Contribution Statement}
K.C.: Methodology feature-level, Software, Formal Analysis, Investigation,
Visualisation, Writing -- Original Draft.
J.P.: Methodology signal-level, Software, Formal Analysis, Investigation,
Visualisation, Writing -- Original Draft.
J.\.Z.: Conceptualisation, Methodology,  Supervision, Resources, Writing -- Original Draft, Review \& Editing.

\section*{Declaration of Competing Interests}
The authors declare that they have no known competing financial
interests or personal relationships that could have appeared to
influence the work reported in this paper.

\section*{Data Availability}
The ELM19 database is available from Elmiko Biosignals sp.~z~o.o.\
upon reasonable request to
\href{mailto:EEG.database@elmiko.pl}{EEG.database@elmiko.pl}.

% -----------------------------------------------------------------------
\appendix
\section{Supplementary Tables}
\label{app:tables}

\begin{table}[htbp]
  \centering
  \caption{Summary statistics of patient age (years) and number of
    normal recordings ($n$) by institution for the $\mathrm{ELM}_n$
    subset. Institutions are ordered by anonymised identifier.}
  \label{tab:age_stats}
  \begin{tabular}{lrrrrrr}
    \toprule
    \textbf{ID} & \textbf{$n$} & \textbf{Mean} & \textbf{Median} &
    \textbf{SD} & \textbf{Min} & \textbf{Max} \\
    \midrule
    H1  &   318 & 46.1 & 45.9 & 15.4 & 18.0 & 70.0 \\
    H2  &   691 & 48.1 & 49.2 & 13.7 & 18.4 & 69.9 \\
    H3  &   473 & 46.3 & 47.9 & 14.7 & 18.0 & 70.0 \\
    H4  &   237 & 26.1 & 21.7 & 10.8 & 16.0 & 69.7 \\
    H5  &   315 & 37.6 & 36.1 & 13.3 & 16.0 & 70.0 \\
    H6  & 1\,148 & 39.0 & 36.9 & 12.6 & 16.1 & 69.7 \\
    H7  &   752 & 47.6 & 50.8 & 16.5 & 16.0 & 69.9 \\
    H8  & 1\,033 & 43.5 & 43.6 & 16.0 & 16.0 & 70.0 \\
    H9  &   336 & 35.0 & 32.3 & 13.4 & 16.0 & 69.1 \\
    H10 & 6\,222 & 37.7 & 35.8 & 15.4 & 16.0 & 70.0 \\
    H11 &   166 & 45.6 & 46.6 & 15.6 & 18.0 & 70.0 \\
    H12 &   232 & 29.4 & 22.4 & 14.7 & 16.0 & 69.9 \\
    H13 &   110 & 44.7 & 44.1 & 16.6 & 18.2 & 69.9 \\
    H14 & 2\,445 & 41.5 & 41.3 & 15.5 & 16.0 & 69.9 \\
    H15 & 2\,289 & 34.9 & 33.3 & 12.5 & 16.0 & 69.8 \\
    H16 & 2\,553 & 45.6 & 46.5 & 14.8 & 16.0 & 70.0 \\
    H17 &   889 & 44.2 & 45.4 & 14.5 & 16.3 & 70.0 \\
    H18 &   488 & 42.8 & 43.0 & 16.6 & 16.1 & 69.9 \\
    H19 &   101 & 48.6 & 52.1 & 15.4 & 18.7 & 69.4 \\
    H20 &   363 & 43.5 & 43.9 & 14.2 & 16.3 & 70.0 \\
    H21 &    90 & 28.4 & 19.5 & 15.1 & 16.0 & 66.3 \\
    H22 & 1\,271 & 45.5 & 46.0 & 15.1 & 16.0 & 70.0 \\
    H23 &    86 & 44.2 & 46.5 & 15.5 & 16.4 & 69.8 \\
    H24 &   783 & 42.3 & 43.3 & 15.4 & 16.0 & 70.0 \\
    H25 &   983 & 38.2 & 37.3 & 15.0 & 16.1 & 70.0 \\
    H26 &   704 & 41.0 & 41.0 & 16.0 & 16.0 & 70.0 \\
    H27 &   307 & 41.3 & 40.9 & 13.8 & 16.0 & 70.0 \\
    H28 &   410 & 46.2 & 48.2 & 15.6 & 17.3 & 69.9 \\
    H29 &   677 & 26.3 & 17.8 & 15.0 & 16.0 & 69.8 \\
    H30 &    96 & 20.3 & 17.6 &  8.3 & 16.0 & 68.3 \\
    \bottomrule
  \end{tabular}
\end{table}

\begin{table}[htbp]
  \centering
  \caption{Acquisition metadata for the 30 institutions in
    $\mathrm{ELM}_s$. Hardware model (A or B), software version,
    sampling frequency, and acquisition-level filter settings
    (high-pass and low-pass, in Hz) are extracted from EDF file
    headers; dashes indicate missing documentation. Where multiple
    values appear, different recordings at the same site used
    different settings.}
  \label{tab:acquisition_metadata}
  \small
  \begin{tabular}{lcllll}
    \toprule
    \textbf{ID} & \textbf{HW} &
    \textbf{Software} & \textbf{Sfreq} &
    \textbf{HP} & \textbf{LP} \\
    \midrule
    H1  & ---  & ---                 & 250       & 0          & 125 \\
    H2  & B    & 12.2 / 14.1        & 250, 256  & 0, 0.16   & 104, 125 \\
    H3  & ---  & ---                 & 250       & 0          & 125 \\
    H4  & B    & 14.1.0.1           & 250, 256  & 0, 0.16   & 104, 125 \\
    H5  & ---  & ---                 & 250       & 0          & 125 \\
    H6  & ---  & ---                 & 250       & 0          & 125 \\
    H7  & B    & 10.3.8.1           & 250, 256  & 0, 0.16   & 125, 1500 \\
    H8  & A    & 10.3.8.1           & 250       & 0, 0.16   & 97, 125 \\
    H9  & B    & 12                  & 256       & 0.16       & 104, 1500 \\
    H10 & A    & 6                   & 250       & 0, 0.16   & 47, 125 \\
    H11 & A    & 6.82                & 250       & 0, 0.16   & 71, 125 \\
    H12 & ---  & ---                 & 250, 256, 600 & 0, 0.16 & 97, 104, 125, 300, 1500 \\
    H13 & A    & 10.3 / 15.0        & 250       & 0, 0.16   & 47, 97, 125 \\
    H14 & ---  & ---                 & 250       & 0          & 125 \\
    H15 & A    & 6.90.2.0            & 250       & 0          & 125 \\
    H16 & ---  & ---                 & 250       & 0.16       & 97 \\
    H17 & A    & 6.90.3.1            & 250       & 0          & 125 \\
    H18 & A    & 8.0.7.8             & 250       & 0          & 125 \\
    H19 & A    & 6.90.3.1            & 250       & 0          & 125 \\
    H20 & A    & 10.3.1.3            & 250       & 0, 0.16   & 97, 125 \\
    H21 & ---  & ---                 & 250       & 0, 0.16   & 97, 125 \\
    H22 & A    & 6.90.3.1            & 250       & 0          & 125 \\
    H23 & A    & 6.90.0.8            & 250       & 0          & 125 \\
    H24 & A    & ---                 & 250, 300  & 0          & 125, 150 \\
    H25 & A    & ---                 & 250       & 0          & 125 \\
    H26 & A    & 6.90 / 10.3 / 14.0 & 250       & 0, 0.16   & 47, 97, 125 \\
    H27 & A    & 6.82                & 250, 600  & 0          & 125, 300 \\
    H28 & A    & 6.61                & 250       & 0          & 125 \\
    H29 & ---  & ---                 & 250, 256  & 0, 0.16   & 47, 97, 104, 125 \\
    H30 & A    & 10.3.8.1            & 250       & 0, 0.16   & 97, 125 \\
    \bottomrule
  \end{tabular}
\end{table}

\newpage

\section{Supplementary Figures}
\label{app:Figs}
\begin{figure}[htbp]
  \centering
  \includegraphics[width=\textwidth]{Figs/raw_amplitude_by_hospital.png}
  \caption{Raw signal amplitude measured prior to preprocessing (median of per-channel MAD, $\mu$V)
    by institution for normal recordings in $\mathrm{ELM}_n$.
    The $y$-axis is log-scaled.
    The two-order-of-magnitude range reflects differences in amplifier
    gain settings across sites.}
  \label{fig:raw_amplitude}
\end{figure}

\begin{figure}[htbp]
  \centering
  \includegraphics[width=0.7\textwidth]{manuscript_Biomedical_Signal_Processing_and_Control/Figs/dom_mcc_newlabels.png}
  \caption{Residual site effect across the signal-level training main
    variants - DANN and MTL, quantified by domain MCC evaluated on the source pool.
    MINet-MTL (no gradient reversal, final epoch) retains strong
    site-discriminative information while MINet-DANN (final and best epoch) reduces domain MCC to
    near-zero values, confirming successful adversarial harmonisation.}
  \label{fig:dann_domain_mcc}
\end{figure}


\begin{figure}[htbp]
  \centering
  \includegraphics[width=\textwidth]{manuscript_Biomedical_Signal_Processing_and_Control/Figs/mcc_new_labels_h1-30.png}
  \caption{Per-institution neuroscreening MCC under the LOSO protocol
    for MINet (unharmonised baseline), MINet-DANN, MINet-MTL and MINet-DANN-CF. Performance is reported for the optimal model checkpoint (Best epoch) -- the same as for the AUC plot.}
  \label{fig:all_strategies_mcc}
\end{figure}
% -----------------------------------------------------------------------
\newpage
\bibliographystyle{elsarticle-num}
\bibliography{bibliography}

\end{document}